% concatenated from frame.tex
\documentclass[reqno,11pt]{amsart}

\usepackage{amsmath,amssymb,amsfonts,amsthm,mathtools}
\usepackage{mathrsfs}
\usepackage{bm}
\usepackage{enumitem}
\usepackage{array,tabularx}
\usepackage{microtype}
\usepackage[colorlinks=true,citecolor=blue,linkcolor=blue,urlcolor=blue]{hyperref}
\usepackage[nameinlink,capitalise,noabbrev]{cleveref}
\allowdisplaybreaks

\setlength{\textwidth}{6.35in}
\setlength{\textheight}{9.1in}
\setlength{\oddsidemargin}{0.05in}
\setlength{\evensidemargin}{0.05in}
\setlength{\topmargin}{-0.25in}

\newtheorem{theorem}{Theorem}[section]
\newtheorem{proposition}[theorem]{Proposition}
\newtheorem{lemma}[theorem]{Lemma}
\newtheorem{corollary}[theorem]{Corollary}
\newtheorem{conjecture}[theorem]{Conjecture}
\theoremstyle{definition}
\newtheorem{definition}[theorem]{Definition}
\newtheorem{remark}[theorem]{Remark}
\newtheorem{notation}[theorem]{Notation}

\newcommand{\SM}{SM}
\newcommand{\FM}{FM}
\newcommand{\cN}{\mathcal N}
\newcommand{\Rop}{\mathscr R}
\newcommand{\cA}{\mathcal A}
\newcommand{\cB}{\mathcal B}
\newcommand{\cQ}{\mathcal Q}
\newcommand{\cC}{\mathcal C}
\newcommand{\cE}{\mathcal E}
\newcommand{\Sym}{\operatorname{Sym}}
\newcommand{\End}{\operatorname{End}}
\newcommand{\SO}{\operatorname{SO}}
\newcommand{\Spin}{\operatorname{Spin}}
\newcommand{\SU}{\operatorname{SU}}
\newcommand{\tr}{\operatorname{tr}}
\newcommand{\Id}{\operatorname{Id}}
\newcommand{\vol}{\operatorname{vol}}
\newcommand{\secg}{\operatorname{sec}_{g}}
\newcommand{\tf}{\operatorname{tf}}
\newcommand{\bbR}{\mathbb R}
\newcommand{\bbC}{\mathbb C}
\newcommand{\cD}{\mathcal D}
\newcommand{\cP}{\mathcal P}
\newcommand{\cS}{\mathcal S}
\newcommand{\cU}{\mathcal U}
\newcommand{\Harm}{\mathscr H}
\newcommand{\Hook}{\mathbb S_{(3,1)}}
\newcommand{\inner}[2]{\left\langle #1,#2\right\rangle}
\newcommand{\norm}[1]{\left\lVert #1\right\rVert}

\title[Frame Flows]
{Pinching and Tensorial Rigidity for Ergodicity of Frame Flows}

\author{Heng Zhang}
\address{School of Mathematical Sciences, University of Science and Technology of China, Hefei, China}
\email{hengz@mail.ustc.edu.cn}

\author{Shuhao Zhang}
\address{School of Mathematical Sciences, University of Science and Technology of China, Hefei, China}
\email{yichen12@mail.ustc.edu.cn}

\subjclass[2020]{37A25, 37D40, 53C20, 53C21, 53C24}
\keywords{frame flow, negatively curved manifolds, Pestov identity, Weitzenb\"ock formula, Killing tensor}

\begin{document}

\begin{abstract}
Let $(M^n,g)$ denote a closed oriented negatively curved Riemannian manifold. For such manifolds, the oriented frame flow is known to be ergodic for all odd dimensions $n\neq7$. We prove Brin's quarter-pinching conjecture when
$n=4$, and when $n\equiv2\pmod4$ with $n\neq134$: strict $1/4$-pinching
implies ergodicity of the oriented frame flow. We also prove that if
$n\equiv0\pmod4$ and $n\geq12$, then $5/13$-pinching implies
ergodicity. In the exceptional dimensions $7$, $8$, and $134$, we prove
ergodicity under strict $0.4661...$-, $0.5358...$-, and
$2/5$-pinching, respectively.
These results substantially improve the
corresponding bounds obtained by Ceki\'c--Lefeuvre--Moroianu--Semmelmann. 
\end{abstract}

\maketitle


\section{Introduction and main results}\label{sec:introduction}

Let $(M^n,g)$, $n\geq3$, be a smooth closed oriented Riemannian manifold with
negative sectional curvature. We write
$$
  \pi:\SM\longrightarrow M,
  \qquad
  \SM=\{(x,v):x\in M,\ v\in T_xM,\ |v|_g=1\}
$$
for its unit tangent bundle, and denote by $\FM\to M$ the oriented
orthonormal frame bundle.  A point of $\FM$ over $x$ will be represented by
an oriented orthonormal basis
$
  (v,e_2,\ldots,e_n)
$
of $T_xM$.  Retaining only the first vector defines a principal bundle
$$
  p:\FM\longrightarrow\SM,
  \qquad
  p(v,e_2,\ldots,e_n)=(x,v),
$$
with structure group $\SO(n-1)$.

Let $(\varphi_t)_{t\in\bbR}$ be the geodesic flow on $\SM$.  The oriented
frame flow $(\Phi_t)_{t\in\bbR}$ is obtained by moving $(x,v)$ under
$\varphi_t$ and parallel-transporting the remaining vectors along the same
geodesic.  Thus
$
  p\circ\Phi_t=\varphi_t\circ p.
$
The flow preserves the smooth probability measure formed from Liouville
measure on $\SM$ and Haar measure on the $\SO(n-1)$-fibres, and ergodicity
will always refer to this measure.  Because negative curvature makes the
geodesic flow Anosov, the frame flow is a compact isometric extension of an
Anosov flow and is one of the basic examples of partial hyperbolicity; see
\cite{HasselblattPesin06}.

For $\delta\in(0,1]$, we say that $(M,g)$ is \emph{$\delta$-pinched} if there
exists $\kappa_0>0$ such that
\begin{equation}\label{eq:delta-pinching-intro}
  -\kappa_0\leq\secg(\sigma)\leq-\delta\kappa_0
  \qquad\text{for every two-plane }\sigma\subset TM.
\end{equation}
It is \emph{strictly $\delta$-pinched} when the right-hand inequality is
strict for every $\sigma$.  On a closed manifold, compactness of the
Grassmann bundle shows that strict $\delta_0$-pinching is equivalent to
$\delta$-pinching for some $\delta>\delta_0$.  In particular, the strict
quarter-pinching hypothesis can be expressed by choosing a pinching
constant $\delta>1/4$ in \eqref{eq:delta-pinching-intro}.

The value $1/4$ is forced by geometry rather than by the known estimates.
Indeed, a parallel complex structure produces a proper invariant reduction
of the oriented frame bundle, so a negatively curved K\"ahler manifold does
not have an ergodic full frame flow.  Berger's pinching theorem places such
metrics at or below the quarter-pinched threshold, and complex hyperbolic
metrics lie on the weak boundary \cite{Berger60}.  This led Brin to the
following conjecture \cite[Conjecture~2.6]{Brin82}.

\begin{conjecture}[Brin]\label{conj:Brin}
If a closed negatively curved Riemannian manifold is strictly
$1/4$-pinched, then its oriented frame flow is ergodic.
\end{conjecture}

Brin's original formulation is stronger: under the same curvature
assumption he predicts that the frame flow is Bernoulli, and he further
conjectures ergodicity and the Bernoulli property whenever the holonomy group
is $\SO(n)$ \cite[Conjectures~2.6 and 2.9]{Brin82}.  The Bernoulli conclusion
is substantially stronger than ergodicity; a recent general mechanism by
which exponential mixing yields Bernoulli behavior is established in
\cite{DolgopyatKanigowskiRodriguezHertz24}.

The unpinched problem is already settled in most odd dimensions.  Brin and
Gromov proved ergodicity for every odd $n\neq7$ without any curvature
pinching beyond negativity \cite{BrinGromov80}.  Brin also showed that
ergodicity holds on an open dense set of negatively curved metrics in the
$C^3$ topology \cite[Section~5]{Brin82}, although a uniform theorem in even
dimensions, and the expected unpinched statement in dimension $7$, are
considerably subtler.  The first quantitative results required curvature
close to constant: Brin--Karcher obtained the threshold
$0.8649\ldots$ in even dimensions other than $8$, while
Burns--Pollicott obtained $0.9805\ldots$ in dimensions $7$ and $8$
\cite{BrinKarcher84,BurnsPollicott03}.

A different framework was introduced by
Ceki\'c--Lefeuvre--Moroianu--Semmelmann \cite{CLMS24}.  In the dimensions
relevant to the present paper, their estimates give the thresholds
$0.2928\ldots$ for $n=4$ and $0.2823\ldots$ for $n=6$.  For
$n\equiv2\pmod4$, $n\geq10$, they obtain an explicit increasing sequence of
bounds beginning with $0.2725\ldots$ in dimension ten and converging to
$0.2779\ldots$.  Their theorem also treats the complementary cases: the
bounds in dimensions divisible by $4$ decrease from $0.5948\ldots$ in
dimension twelve toward $0.5572\ldots$, and the exceptional dimensions
$7$, $8$, and $134$ are covered by the respective thresholds
$0.4962\ldots$, $0.6212\ldots$, and $0.5788\ldots$.  The conceptual
change is as important as the numerical improvement: the earlier arguments
were either essentially topological or based on the geometry of the
universal cover, whereas the method of \cite{CLMS24} converts
non-ergodicity into invariant tensorial data on $\SM$ and then applies a
twisted Pestov identity \cite{BrinGromov80,BrinKarcher84,BurnsPollicott03}.

The ergodicity question is part of a broader theory of compact isometric
extensions of hyperbolic flows.  For frame flows in dimensions $n\geq4$, ergodicity implies mixing
\cite{Lefeuvre23}. Moreover, since the structure group
$\SO(n-1)$ is compact semisimple in this range, ergodicity is
equivalent to rapid mixing by
\cite{CekicLefeuvreSemiclassical24}. Related quantitative criteria,
applying in particular to some frame flows of negatively curved
manifolds, are established in \cite{PollicottZhang25}.  When a K\"ahler
structure obstructs ergodicity of the full frame flow, the corresponding
unitary frame flow becomes the natural object: if the complex dimension is
$m$, Brin--Gromov's argument gives ergodicity for $m$ odd or $m=2$, while
holomorphic-pinching results cover even $m\neq4,28$
\cite{BrinGromov80,CLMSUnitary24}.  There is also an analogous theory for an
arbitrary Euclidean vector bundle $(E,\nabla)\to M$ with an orthogonal
connection: if $\operatorname{rank}E=r\leq\sqrt n$ and the holonomy group is
$\SO(r)$, then the associated frame flow is ergodic
\cite{CekicLefeuvreConnections25}.  This gives further evidence for Brin's
holonomy conjecture and indicates that the methods extend well beyond the
tangent bundle.  For an overview of this circle of ideas, see
\cite{CLMS22Survey}.

Our results divide naturally into two groups.  The first consists of
stable dimensional families, where the obstruction is an invariant normal
vector or an invariant orthogonal projector.  The second consists of the
three exceptional dimensions in which the transitivity-group classification
permits additional form-valued structures.

\begin{theorem}\label{thm:stable-main}
Let $(M^n,g)$ be a closed oriented Riemannian manifold of dimension $n\geq4$
with negative sectional curvature.
\begin{enumerate}[label=\textup{(\roman*)}]
\item If $g$ is strictly $1/4$-pinched and
$$
  n=4,
  \qquad\text{or}\qquad
  n\equiv2\pmod4\quad\text{with}\quad n\neq134,
$$
then the oriented frame flow on $\FM$ is ergodic.
\item If $g$ is $5/13$-pinched and
$$
  n\equiv0\pmod4,
  \qquad n\geq12,
$$
then the oriented frame flow on $\FM$ is ergodic.
\end{enumerate}
\end{theorem}

For the exceptional dimensions, set
$$
\delta_7^\sharp:=\frac{385+36\sqrt6}{927+36\sqrt6}=0.4661\ldots, \qquad
 \delta_8^\sharp:=4-2\sqrt3=0.5358\ldots.
$$

\begin{theorem}\label{thm:exceptional-main}
Let $(M^n,g)$ be closed, oriented, and negatively curved.
\begin{enumerate}[label=\textup{(\roman*)}]
\item If $n=7$ and $g$ is strictly $\delta_7^\sharp$-pinched, then the
oriented frame flow is ergodic.
\item If $n=8$ and $g$ is strictly $\delta_8^\sharp$-pinched, then the
oriented frame flow is ergodic.
\item If $n=134$ and $g$ is strictly $2/5$-pinched, then the oriented frame
flow is ergodic.
\end{enumerate}
\end{theorem}





\begin{remark}\label{rem:intro-improvement}
The pinching constants $5/13$, $\delta_7^\sharp$, $\delta_8^\sharp$, and
$2/5$ may be slightly lowered by modifying our method via some new estimates.  The exact optimization
for Theorem \ref{thm:stable-main}\textup{(ii)} is recorded in
Remark \ref{rem:optimized-constant}; we retain $5/13$ in the statement
because it yields rational coefficients and visibly positive gaps in
every irreducible quartic summand.
\end{remark}


We now outline the proof and the main new ingredients.  Let
\(\cN\to\SM\) denote the normal bundle, whose fibre at \((x,v)\) is
\(v^\perp\subset T_xM\).  The initial dynamical reduction is the one
developed by Brin
\cite{Brin75Topology,Brin75Transitivity}
and refined in
\cite{CekicLefeuvreHolonomy25,CLMS24}.  If the frame flow is not
ergodic, then, after passing to a finite Riemannian cover, its
transitivity group fixes one of the tensorial structures listed in
Theorem~\ref{thm:CLMS-input}.  Non-Abelian Liv\v{s}ic theory realizes
this fixed tensor as a smooth flow-invariant section of the normal
bundle, of its symmetric square, or of an exterior power.

Every such invariant section has finite vertical Fourier content
\cite{GPSU16}.  Its highest Fourier coefficient can be viewed on the
base manifold as a symmetric tensor, twisted by the relevant auxiliary
bundle, which satisfies both a conformal-Killing equation and an
algebraic normality constraint.  The localized twisted Pestov identity
then forces an associated curvature form to be nonpositive.  Thus the
analytic part of the proof consists in proving that the same curvature
form is strictly positive under the pinching assumptions.  A recurring
principle in our argument is that the Jacobi and twisting contributions
must be retained as a single expression: estimating them separately,
as in the coarser bounds, loses precisely the interaction needed near
the quarter-pinched threshold.

We first consider an invariant normal vector field.  A new estimate for
the combined tangent-twist Pestov form excludes all odd Fourier degrees
at least five already at the quarter-pinched endpoint.  Consequently,
only the linear and cubic modes remain.  A purely linear obstruction
would define an orthogonal nearly K\"ahler structure, which is ruled out
by the structure theory of nearly K\"ahler manifolds and Berger's
pinching theorem
\cite{Gray76,Nagy02,Berger60}.  When the two remaining modes are
coupled, they assemble naturally into a generalized Killing tensor of
Young type \((3,1)\).  The resulting differential equation is treated
by Weitzenb\"ock methods.  In dimension $4$, the essential ingredients
are a quantitative Finsler--Thorpe lemma and the
\(\Spin(4)=\SU(2)_+\times\SU(2)_-\) decomposition.  Dimension $6$ is
handled by a direct low-dimensional tangent-twist estimate.  In even
dimensions at least ten, the hook tensor splits into a totally
trace-free component and two trace components.  The trace-free part is
controlled by a new cubic Weitzenb\"ock identity which isolates the
orthogonal action on the first three tensor slots, while the two trace
parts are treated by the ordinary Weitzenb\"ock formula.  Together,
these arguments establish the quarter-pinched result for the
normal-vector branch.

For an invariant orthogonal projector, we similarly retain the full
curvature expression associated with the symmetric-square twist.  A
uniform estimate eliminates all even Fourier degrees at least $6$
under \(5/13\)-pinching.  The quartic mode requires a separate
finite-dimensional argument: its admissible coefficient space has a
multiplicity-free decomposition into three irreducible orthogonal-group
modules, and the relevant invariant quadratic forms can be evaluated
exactly on each summand.  The quadratic mode has additional rigidity.
The projector equation and normality produce a $4$-tensor of Young
type \((2,2)\); the resulting derivative antisymmetries, together with
the first Bianchi identity, give a sharper estimate for the mixed
curvature term.  This excludes the projector branch in dimensions
divisible by $4$, and also deals with the projector alternative in
dimension eight.

The exceptional \(G_2\) and \(E_7\) branches require a different use of
the invariant structure.  Here it is advantageous to keep the entire
finite Fourier expansion rather than only its highest coefficient.
After homogenization, the invariant form becomes a generalized Killing
tensor on the base while retaining both its normality and its fixed
orbit type.  We prove an orbit-type curvature identity which decomposes
the relevant curvature term into three pieces: a radial Jacobi term, a
radial--tangential mixed term, and the trace of the curvature operator
on the orthogonal complement of the stabilizer algebra.  In both
applications this complement is irreducible, so the orbit differential
is a homothety and compact-group averaging produces a tight frame.
For the \(G_2\) orbit, the standard decomposition of two-forms gives the
dimension-eight threshold.  For the \(E_7\) Cartan orbit, a tight frame
generated by commuting decomposable two-forms gives the
dimension-\(134\) threshold stated in
Theorem~\ref{thm:exceptional-main}.

Finally, the $7$-dimensional complex-structure branch is treated by
a separate cubic refinement.  The higher odd Fourier modes are first
removed by the tangent-twist estimate.  For the remaining cubic mode,
the normality relation produces a tangential vector harmonic of degree
two.  A sharp \(L^2\) comparison between the cubic tensor and this
contraction, combined with the lowering term retained in a second
localized Pestov estimate, excludes the cubic mode under the constant
in Theorem~\ref{thm:exceptional-main}\textup{(i)}.  The remaining
linear mode induces a nearly parallel \(G_2\)-structure and is
incompatible with negative sectional curvature.

These vanishing results exclude every obstruction supplied by
Theorem~\ref{thm:CLMS-input}, and therefore prove
Theorems~\ref{thm:stable-main} and
\ref{thm:exceptional-main}.

For comparison, Table~\ref{tab:comparison} summarizes the obstruction supplied by
the dynamical reduction, the bounds obtained in \cite{CLMS24}, and the estimates
proved here.

\begin{table}[!ht]
\centering
\footnotesize
\setlength{\tabcolsep}{3pt}
\renewcommand{\arraystretch}{1.18}
\begin{tabularx}{\textwidth}{@{}
  >{\raggedright\arraybackslash}p{0.82in}
  >{\raggedright\arraybackslash}p{0.92in}
  >{\raggedright\arraybackslash}p{1.08in}
  >{\raggedright\arraybackslash}p{0.82in}
  >{\raggedright\arraybackslash}X@{}}
\hline
Dimension(s) & Obstruction & Threshold in \cite{CLMS24} & Present threshold & Principal new ingredient \\
\hline
$4$ & odd normal vector & $0.2928\ldots$ & strict $1/4$ & quantitative Finsler--Thorpe lemma and the $\Spin(4)$ decomposition \\
$n\equiv2\pmod4$, $n\neq134$ & odd normal vector & $0.2823\ldots$ for $n=6$; $0.2725\ldots\nearrow0.2779\ldots$ for $n\geq10$ & strict $1/4$ & tangent-twist Pestov estimate and $(3,1)$-type Weitzenb\"ock identities \\
$4\mid n$, $n\geq12$ & even orthogonal projector & $0.5948\ldots\searrow0.5572\ldots$ & $5/13$ & symmetric-square Pestov curvature form, admissible quartic module, and quadratic projector symmetry \\
$7$ & orthogonal complex structure on $\cN$ & $0.4962\ldots$ & strict $\delta_7^\sharp=0.4661\ldots$ & cubic contraction and retained lowering term \\
$8$ & $G_2$ form or projector & $0.6212\ldots$ & strict $\delta_8^\sharp=0.5358\ldots$ & orbit-type curvature identity and the projector estimate \\
$134$ & normal vector or $E_7$ Cartan form & $0.5788\ldots$ & strict $2/5$ & $(3,1)$-type Weitzenb\"ock estimate and the orbit-type curvature identity \\
\hline
\end{tabularx}
\caption{Comparison of the dimensional branches and pinching thresholds.}
\label{tab:comparison}
\end{table}

\medskip
\noindent\emph{Organization of the paper.}
Section \ref{sec:preliminaries} records the geometric conventions, the vertical
Fourier decomposition, the localized Pestov identities, and the dynamical
reduction to invariant tensorial obstructions.  Section \ref{sec:normal-vector}
treats the invariant normal-vector obstruction: first the high-degree
reduction, then the generalized Killing tensor of Young type $(3,1)$, and
finally the estimates in dimensions $4$, $6$, and the stable range.  Section
\ref{sec:symmetric-projector} treats the invariant projector obstruction in
Fourier degrees at least $6$, $4$, and $2$.  Section
\ref{sec:exceptional-dimensions} develops the general orbit-type curvature
identity and applies it to the $G_2$ and $E_7$ branches, followed by the
separate $7$-dimensional refinement.  Section \ref{sec:mainproof} combines
these vanishing results and proves Theorems \ref{thm:stable-main} and
\ref{thm:exceptional-main}.  Appendices \ref{app:JM}--\ref{app:projector-optimization}
contain the representation-theoretic and coefficient computations used in
the main text.

\section{Dynamical reduction, conventions, and Pestov identities}\label{sec:preliminaries}

\subsection{Pinching, normalization, and curvature conventions}

In algebraic calculations, $V$ denotes a Euclidean vector space, and we use
its metric to identify $V$ with $V^*$. Parentheses and brackets around
indices denote normalized symmetrization and alternation, respectively;
repeated indices are summed.

We use the Riemann curvature convention
$$
  R^g(u,v)w
  :=\nabla_u\nabla_vw-\nabla_v\nabla_uw-\nabla_{[u,v]}w,
$$
and write
$$
  R^g(u,v,w,z):=\inner{R^g(u,v)w}{z},
  \qquad
  \secg(u\wedge v)
  =\frac{R^g(u,v,v,u)}{|u\wedge v|^2}.
$$
Define the associated \emph{positive curvature operator}
\(
  \Rop^g:\Lambda^2TM\to\Lambda^2TM
\)
by
\begin{equation}\label{eq:Rop-definition}
  \inner{\Rop^g(u\wedge v)}{w\wedge z}
  :=-R^g(u,v,z,w)=R^g(u,v,w,z).
\end{equation}
Thus
\begin{equation*}
  \inner{\Rop^g(u\wedge v)}{u\wedge v}
  =-\secg(u\wedge v)|u\wedge v|^2.
\end{equation*}
For a self-adjoint $\cB\in\End(\Lambda^2V)$, write
$$
  \cB(\xi,\eta):=\inner{\cB\xi}{\eta},
  \qquad
  \cB(u,v,w,z):=\inner{\cB(u\wedge v)}{w\wedge z},
$$
and set
$$
  \sec_{\cB}(u\wedge v)
  :=\frac{\cB(u\wedge v,u\wedge v)}{|u\wedge v|^2},
  \qquad
  \sec_{\cB}(u,v):=\sec_{\cB}(u\wedge v).
$$
Relative to an orthonormal basis $(e_i)$, write
$$
  e_{ij}:=e_i\wedge e_j,
  \qquad
  \cB_{ijkl}:=\cB(e_i,e_j,e_k,e_l).
$$
We write $I$ for the identity endomorphism of $\Lambda^2V$; thus
$\sec_I\equiv1$.

Suppose that $g$ is $\delta$-pinched at the scale $\kappa_0>0$ in the sense
of \eqref{eq:delta-pinching-intro}. We then set
\begin{equation}\label{eq:normalized-curvature}
  \cA:=\kappa_0^{-1}\Rop^g.
\end{equation}
The pinching condition becomes
\begin{equation*}
  \delta\leq\sec_{\cA}\leq1.
\end{equation*}
All algebraic curvature estimates below are stated for $\cA$. Since every
curvature expression used in the proof is linear in the curvature operator,
the positive factor $\kappa_0$ never affects a sign.

Let $E$ be a Euclidean $\SO(n)$-module with infinitesimal representation
\(
  \rho:\mathfrak{so}(n)\to\mathfrak{so}(E).
\)
For any self-adjoint $\cB\in\End(\Lambda^2\bbR^n)$ and any orthonormal basis
$(E_\alpha)$ of $\Lambda^2\bbR^n\simeq\mathfrak{so}(n)$, define
\begin{equation*}
  \cQ_E(\cB)
  :=\sum_{\alpha,\beta}
  \inner{\cB E_\alpha}{E_\beta}\,
  \rho(E_\alpha)^*\rho(E_\beta).
\end{equation*}
This is independent of the chosen basis. If $\cB\succeq0$ as an operator on
$\Lambda^2\bbR^n$, then
\begin{equation}\label{eq:Qpositive}
  \cQ_E(\cB)\succeq0.
\end{equation}
Indeed, after diagonalizing
$\cB=\sum_\alpha\lambda_\alpha E_\alpha\otimes E_\alpha$, one has
$$
  \cQ_E(\cB)
  =\sum_\alpha\lambda_\alpha
  \rho(E_\alpha)^*\rho(E_\alpha).
$$

Let $q_E(R^g)$ be the standard Weitzenb\"ock curvature endomorphism on the
bundle associated with $E$. With the conventions above,
\begin{equation}\label{eq:q-Q-exact-convention}
  q_E(R^g)=-\cQ_E(\Rop^g)
  =-\kappa_0\cQ_E(\cA).
\end{equation}
This is the convention of \cite[Section~3]{CLMSW25}.

\subsection{The unit tangent bundle and the normal bundle}

The geodesic vector field on $\SM$ is denoted by $X$, and its flow is
$$
  \varphi_t(x,v)
  =(\gamma_{x,v}(t),\dot\gamma_{x,v}(t)).
$$

\begin{definition}
The normal bundle $\cN\to\SM$ is
$$
  \cN_{(x,v)}:=v^\perp\subset T_xM.
$$
We use $X$ both for the geodesic vector field and for the induced covariant
derivative on pullback tensor bundles over $\SM$.
A section $f\in C^\infty(\SM,\cN)$ is \emph{flow-invariant} if $Xf=0$.
\end{definition}

Thus $Xf=0$ means that $f(\varphi_t(x,v))$ is obtained from $f(x,v)$ by
parallel transport along $\gamma_{x,v}$.

\subsection{Vertical spherical harmonics}

For each $x\in M$, the fibre $S_xM$ is a round $(n-1)$-sphere. Let
$\Omega_k\to M$ be the vector bundle of spherical harmonics of degree $k$ on
these fibres. Every smooth section
$f\in C^\infty(\SM,\pi^*TM)$ has a vertical Fourier expansion
$$
  f=\sum_{k\geq0}f_k,
  \qquad
  f_k\in C^\infty(M,\Omega_k\otimes TM),
$$
with convergence in the smooth topology. The section is odd if
$f(x,-v)=-f(x,v)$, equivalently if only odd values of $k$ occur.

The geodesic vector field splits into raising and lowering operators
$$
  X=X_++X_- ,
  \qquad
  X_+:\Omega_k\otimes TM\longrightarrow\Omega_{k+1}\otimes TM,
  \qquad
  X_-:\Omega_k\otimes TM\longrightarrow\Omega_{k-1}\otimes TM.
$$
If $f$ has finite Fourier degree $k$ and $Xf=0$, then its top coefficient
satisfies
$$
  X_+f_k=0.
$$
Via the standard identification between spherical harmonics and trace-free
symmetric tensors, there is a unique
$$
  K_k\in C^\infty(M,\Sym^k_0T^*M\otimes TM)
$$
whose evaluation on $v^{\otimes k}$ is $f_k(x,v)$. The equation $X_+f_k=0$
is the twisted conformal Killing equation. We refer to
\cite[Sections~2.2 and 4.1]{CLMS24} and \cite{CLMSW25} for the precise
normalizations.  Further background on Killing and conformal Killing symmetric
tensors is given in
\cite{DairbekovSharafutdinov10,HeilMoroianuSemmelmann16}.

\subsection{Normalized localized Pestov identities}
\label{sec:normalized-pestov}

We collect the normalizations and sign conventions for the Pestov identities
used throughout the paper.  Let $(E,\nabla^E)\to M$ be a Euclidean vector
bundle with an orthogonal connection, and equip $\SM$ with the Liouville
measure
$$
  d\mu(x,v)=d\vol_g(x)\,d\sigma_x(v),
  \qquad \int_{S_xM}d\sigma_x=1.
$$
All adjoints and norms in this subsection are taken with respect to this
measure and the bundle metrics.  The geodesic derivative is skew-adjoint,
$$
  X^*=-X,
$$
and, on the vertical harmonic decomposition,
\begin{equation*}
  \left(X_+\big|_{\Omega_k\otimes E}\right)^*
  =-X_-\big|_{\Omega_{k+1}\otimes E}.
\end{equation*}
In particular, $X_-$ vanishes on degree zero.

For $u\in C^\infty(M,\Omega_k\otimes E)$, our vertical Laplacian is the
nonnegative operator
\begin{equation*}
  \Delta_V^E:=(\nabla_V^E)^*\nabla_V^E,
  \qquad
  \Delta_V^E u=\lambda_k u,
  \qquad
  \lambda_k=k(k+n-2).
\end{equation*}
At $(x,v)\in\SM$, choose an orthonormal basis
$e_1,\ldots,e_{n-1}$ of $v^\perp$.  The vertical derivative is normalized by
\begin{equation*}
  D_a^Vu
  :=\left.\frac{D}{dt}\right|_{t=0}
  u\bigl(x,(\cos t)v+(\sin t)e_a\bigr),
  \qquad
  \nabla_V^Eu=\sum_{a=1}^{n-1}e_a\otimes D_a^Vu,
\end{equation*}
and hence
$|\nabla_V^Eu|^2=\sum_a|D_a^Vu|^2$.

Define the Jacobi and twisting curvature operators by
\begin{equation*}
  \mathbf R_{(x,v)}(w\otimes\xi)
  :=R^g(w,v)v\otimes\xi,
  \qquad
  F^E_{(x,v)}\xi
  :=\sum_{a=1}^{n-1}e_a\otimes R^E(v,e_a)\xi.
\end{equation*}
For $k\geq0$, set
\begin{equation*}
  A_{n,k}:=\frac{(n+k-2)(n+2k-4)}{n+k-3},
  \qquad
  B_{n,k}:=\frac{k(n+2k)}{k+1}.
\end{equation*}
The normalized localized twisted Pestov identity is
\begin{equation}\label{eq:localized-pestov-identity}
\begin{split}
  A_{n,k}\|X_-u\|^2-B_{n,k}\|X_+u\|^2+\|Z_k(u)\|^2
  ={}&\inner{\mathbf R\nabla_V^Eu}{\nabla_V^Eu}_{L^2}+\inner{F^Eu}{\nabla_V^Eu}_{L^2}.
\end{split}
\end{equation}
Here $Z_k(u)$ is characterized by the decomposition
\begin{equation*}
  \nabla_H^Eu
  =\frac1{k+1}\nabla_V^EX_+u
   -\frac1{n+k-3}\nabla_V^EX_-u+Z_k(u),
  \qquad \operatorname{div}_V^E Z_k(u)=0.
\end{equation*}
We use only the nonnegativity of $\|Z_k(u)\|^2$.  Formula
\eqref{eq:localized-pestov-identity} is
\cite[Lemma~2.3]{CLMS24}; see also
\cite[Proposition~3.5]{GPSU16} and
\cite[Proposition~6.2]{CLMSW25}.

We now rewrite its curvature side in the positive curvature operator
convention of \eqref{eq:Rop-definition}.  Identify
$\Lambda^2T_xM\simeq\mathfrak{so}(T_xM)$ by
\begin{equation}\label{eq:wedge-skew-convention}
  (a\wedge b)c=\inner{a}{c}b-\inner{b}{c}a.
\end{equation}
For every bundle associated to an orthogonal representation $\rho_E$, one has
\begin{equation}\label{eq:pestov-sign-conversion}
  R^E(v,e_a)=\rho_E\bigl(\Rop^g(v\wedge e_a)\bigr),
  \qquad
  \inner{R^g(e_a,v)v}{e_b}
  =-\Rop^g(v\wedge e_a,v\wedge e_b).
\end{equation}
For a self-adjoint $\cB\in\End(\Lambda^2T_xM)$, define the sign-normalized
Pestov curvature form
\begin{equation}\label{eq:sign-normalized-pestov-form}
\begin{split}
  \mathfrak P_{\cB}^E(u):=\int_{\SM}\bigg[{}
  &\sum_{a,b}\cB(v\wedge e_a,v\wedge e_b)
       \inner{D_a^Vu}{D_b^Vu}-\sum_a\inner{\rho_E(\cB(v\wedge e_a))u}{D_a^Vu}
  \bigg]d\mu.
\end{split}
\end{equation}
The expression is independent of the orthonormal basis of $v^\perp$.
Equations \eqref{eq:localized-pestov-identity} and
\eqref{eq:pestov-sign-conversion} give the single sign-normalized formula
\begin{equation}\label{eq:localized-pestov-positive}
  A_{n,k}\|X_-u\|^2-B_{n,k}\|X_+u\|^2+\|Z_k(u)\|^2
  =-\mathfrak P_{\Rop^g}^E(u).
\end{equation}

For the tangent-bundle twist, let $(e_r)_{r=0}^{n-1}$ be a fixed ambient
orthonormal basis, held fixed during spherical differentiation, and put
$z_r:=\nabla^S\inner{u}{e_r}$.  Since the infinitesimal action on $TM$ is the
standard one, \eqref{eq:sign-normalized-pestov-form} becomes
\begin{equation}\label{eq:tangent-pestov-form}
  \mathfrak P_{\cB}^{TM}(u)
  =\int_{\SM}\sum_{r=0}^{n-1}
  \left\{
    \cB(v\wedge z_r,v\wedge z_r)
    -\cB(u\wedge e_r,v\wedge z_r)
  \right\}d\mu.
\end{equation}
Thus the tangent-twist localized identity is
\begin{equation*}
  A_{n,k}\|X_-u\|^2-B_{n,k}\|X_+u\|^2+\|Z_k(u)\|^2
  =-\mathfrak P_{\Rop^g}^{TM}(u).
\end{equation*}

For the symmetric-square twist, write
$S_{rs}=\inner{S e_s}{e_r}$ and $z_{rs}=\nabla^SS_{rs}$.  The infinitesimal
action is $\rho_{\Sym^2}(A)S=[A,S]$, and the two equal matrix contributions
produce the coefficient two.  Hence
\begin{equation}\label{eq:sym2-pestov-form-normalized}
  \mathfrak P_{\cB}^{\Sym^2}(S)
  =\int_{\SM}\sum_{r,s=0}^{n-1}
  \left\{
    \cB(v\wedge z_{rs},v\wedge z_{rs})
    -2\cB(S(v)e_s\wedge e_r,v\wedge z_{rs})
  \right\}d\mu,
\end{equation}
and the $\Sym^2$-twisted identity is
\begin{equation*}
  A_{n,k}\|X_-S\|^2-B_{n,k}\|X_+S\|^2+\|Z_k(S)\|^2
  =-\mathfrak P_{\Rop^g}^{\Sym^2}(S).
\end{equation*}

For the $\Lambda^p$-twist, let $\rho_{\Lambda^p}$ be the induced infinitesimal
action and write
$J_{\cB}(e_a,e_b)=\cB(v\wedge e_a,v\wedge e_b)$.  Then
\begin{equation*}
\begin{split}
  \mathfrak P_{\cB}^{\Lambda^p}(u)=\int_{\SM}\bigg[{}
  &\sum_{a,b}J_{\cB}(e_a,e_b)
       \inner{D_a^Vu}{D_b^Vu}-\sum_a\inner{
       \rho_{\Lambda^p}(\cB(v\wedge e_a))u}{D_a^Vu}
  \bigg]d\mu,
\end{split}
\end{equation*}
so the $\Lambda^p$-twisted identity is
\begin{equation*}
  A_{n,k}\|X_-u\|^2-B_{n,k}\|X_+u\|^2+\|Z_k(u)\|^2
  =-\mathfrak P_{\Rop^g}^{\Lambda^p}(u).
\end{equation*}

There are two sign consequences used below.  First, if a finite Fourier
expansion $u=\sum_{j=0}^ku_j$ satisfies $Xu=0$, then
$X_+u_k=0$.  Applying \eqref{eq:localized-pestov-positive} to the top
coefficient gives
\begin{equation}\label{eq:top-pestov-sign}
  \mathfrak P_{\Rop^g}^E(u_k)
  =-A_{n,k}\|X_-u_k\|^2-\|Z_k(u_k)\|^2\leq0.
\end{equation}
Second, the full twisted Pestov identity, in the same conventions, is
\cite[Proposition~3.3]{GPSU16}
\begin{equation*}
\begin{split}
  \|X\nabla_V^Eu\|^2-\|\nabla_V^EXu\|^2+(n-1)\|Xu\|^2
  ={}&\inner{\mathbf R\nabla_V^Eu}{\nabla_V^Eu}_{L^2}+\inner{F^Eu}{\nabla_V^Eu}_{L^2}
  =-\mathfrak P_{\Rop^g}^E(u).
\end{split}
\end{equation*}
Consequently, an invariant section itself satisfies
\begin{equation}\label{eq:full-pestov-sign}
  Xu=0
  \quad\Longrightarrow\quad
  \mathfrak P_{\Rop^g}^E(u)
  =-\|X\nabla_V^Eu\|^2\leq0.
\end{equation}
Thus every later nonpositivity statement is obtained directly from either
\eqref{eq:top-pestov-sign} or \eqref{eq:full-pestov-sign}.
\subsection{Dynamical result and the invariant obstructions}

We first record a dynamical result.  For completeness, recall that if
$n=(2a+1)2^b$ and $b=c+4d$ with $0\leq c\leq3$, then the
Radon--Hurwitz number is $\rho(n)=2^c+8d$.

\begin{theorem}
\label{thm:CLMS-input}
Let $(M^n,g)$ be closed, oriented, negatively curved, and suppose that its
oriented frame flow is not ergodic. Then there exists a finite Riemannian
cover $(\widehat M,\widehat g)\to(M,g)$ on which one of the following holds.
\begin{enumerate}[label=\textup{(\roman*)}]
\item If
$$
 n=4,
 \qquad\text{or}\qquad
 n\equiv2\pmod4,\quad n\neq134,
$$
there is an odd unit section
$$
  f\in C^\infty(S\widehat M,\cN),
  \qquad Xf=0,
  \qquad \inner{f(x,v)}v=0.
$$
\item If $n\equiv0\pmod4$ and $n\geq12$, there is a nontrivial even
flow-invariant orthogonal projector
$$
  P\in C^\infty(S\widehat M,\Sym^2\cN),
  \qquad XP=0,
$$
of rank
$$
  1\leq r\leq
  \min\left\{\rho(n)-1,\frac{n-2}{2}\right\}.
$$
\item If $n=7$, there is an odd flow-invariant section
$$
  f\in C^\infty(S\widehat M,\Lambda^2\cN),
  \qquad Xf=0,
$$
whose values are the two-forms of orthogonal complex structures on the
normal $6$-plane.
\item If $n=8$, then either there is an odd flow-invariant $G_2$-structure
$$
  \phi\in C^\infty(S\widehat M,\Lambda^3\cN),
  \qquad X\phi=0,
$$
or there is a nontrivial even flow-invariant orthogonal projector
$$
  P\in C^\infty(S\widehat M,\Sym^2\cN),
  \qquad XP=0,
  \qquad 1\leq\operatorname{rank}P\leq3.
$$
\item If $n=134$, then either there is an odd invariant unit section
$$
  f\in C^\infty(S\widehat M,\cN),
  \qquad Xf=0,
$$
or there is a nonzero odd flow-invariant Lie bracket
$$
  \varphi\in C^\infty(S\widehat M,\Lambda^3\cN),
  \qquad X\varphi=0,
$$
whose values are equivalent to the Cartan three-form of the compact Lie
algebra $\mathfrak e_7$ and whose Fourier degree is at least three.
\end{enumerate}
In all cases the invariant object has finite vertical Fourier degree.
\end{theorem}

\begin{proof}
The invariant structures, their parity, the rank restrictions, and the
degree-three lower bound in the $E_7$ branch are supplied by
\cite[Theorem~3.8 and Lemmas~3.12--3.13]{CLMS24}.  Finite Fourier degree
follows from \cite[Theorem~4.1]{GPSU16}.  The pinching constant and scale
pull back unchanged to the finite cover.
\end{proof}

We shall repeatedly use the following four-index estimate of Berger–Bourguignon–Karcher.
\begin{lemma}\label{lem:BBK}
Let $\cB$ be an algebraic curvature operator satisfying
$$
  a\leq\sec_{\cB}\leq b.
$$
Set $\widehat{\cB}:=\cB-\frac{a+b}{2}I$.  Then, for arbitrary vectors
$u,v,w,z$,
\begin{equation}\label{eq:BBK-vectors}
  |\widehat{\cB}(u,v,w,z)|
  \leq\frac23(b-a)|u|\,|v|\,|w|\,|z|.
\end{equation}
\end{lemma}


\begin{proof}
The unit-vector estimate is \cite[Lemma~3.7]{BourguignonKarcher78};
\eqref{eq:BBK-vectors} follows by homogeneity.  
\end{proof}
Moreover, for decomposable two-forms $\xi$ and $\eta$,
$$
  |\widehat{\cB}(\xi,\eta)|
  \leq\frac23(b-a)|\xi|\,|\eta|.
$$
If in addition $\xi\perp\eta$, then $I(\xi,\eta)=0$ and hence
\begin{equation}\label{eq:BK-simple-form}
  |\cB(\xi,\eta)|
  \leq\frac23(b-a)|\xi|\,|\eta|.
\end{equation}

\section{The invariant normal-vector obstruction}\label{sec:normal-vector}

This section turns the low Fourier obstruction into a generalized Killing
tensor of Young type $(3,1)$ and establishes the two Weitzenb\"ock identities
used in the proof.

\subsection{High-degree reduction}\label{sec:reduction}

\subsubsection{The tangent-twist Pestov curvature form}

Let $V$ be an $n$-dimensional Euclidean space, and choose $c_{n,k}>0$ so that
$$
  |K|^2
  =c_{n,k}\int_{S(V)}|K(x,\ldots,x)|^2\,d\sigma(x)
$$
for $K\in\Sym^k_0V^*\otimes V$. Put
$$
  F(x):=K(x,\ldots,x),
  \qquad
  F_r(x):=\inner{F(x)}{e_r}.
$$
For an algebraic curvature operator $\cB$ on $V$, define
\begin{equation}\label{eq:tangent-twist-form}
  W_{\cB}(K)
  :=c_{n,k}\int_{S(V)}\sum_{r=0}^{n-1}
  \bigl\{\cB(\beta_r,\beta_r)-\cB(\alpha_r,\beta_r)\bigr\}\,d\sigma(x),
  \qquad
  \alpha_r=F\wedge e_r,
  \quad
  \beta_r=x\wedge\nabla^S F_r.
\end{equation}
Here the vectors $(e_r)$ are fixed as ambient constant vectors when the
spherical derivative is taken. The sum is independent of the fixed
orthonormal basis, so at a fixed $x$ one may choose a basis adapted to $x$
and $F(x)$.

Under the tangent-bundle specialization
\eqref{eq:tangent-pestov-form}, the expression
\eqref{eq:tangent-twist-form} is precisely the tangent-twist Pestov curvature form
for the harmonic tensor $K$, with the harmless factor $c_{n,k}>0$ fixing the
harmonic--tensor isometry.  Therefore \eqref{eq:top-pestov-sign} gives, whenever
$K$ is the top Fourier coefficient of an invariant normal section,
\begin{equation*}
  \int_M W_{\Rop^g}(K)\,d\vol_g\leq0.
\end{equation*}
If $g$ is $\delta$-pinched at scale $\kappa_0$ and $\cA$ is defined by
\eqref{eq:normalized-curvature}, linearity gives the equivalent inequality
\begin{equation}\label{eq:tangent-twist-pestov-normalized}
  \int_M W_{\cA}(K)\,d\vol_g\leq0.
\end{equation}

Put
$$
  \phi(x):=\inner{F(x)}x,
  \qquad
  F^\perp(x):=F(x)-\phi(x)x,
$$
and
\begin{equation*}
  N:=c_{n,k}\int|F|^2,
  \qquad
  N_{\parallel}:=c_{n,k}\int \phi^2,
  \qquad
  N_{\perp}:=c_{n,k}\int|F^\perp|^2=N-N_{\parallel}.
\end{equation*}
Unless otherwise indicated, all integrals in this fibrewise calculation are
over $S(V)$.

We estimate the residual for the tangent-twist Pestov form.
\begin{lemma}\label{lem:tangent-twist-residual}
If $0\leq\sec_{\cE}\leq1$, then
\begin{equation*}
  W_{\cE}(K)
  \geq-\frac14\bigl(N_{\perp}+(n-1)N_{\parallel}\bigr)
  -\frac23\sqrt{(n-2)N_{\perp}\,k(k+n-2)N}.
\end{equation*}
\end{lemma}

\begin{proof}
Fix $x\in S(V)$ and write
$$
  F=\phi x+F^\perp,
  \qquad
  F^\perp\perp x.
$$
Since the sum in the definition of $W_{\cE}(K)$ is independent of the
fixed ambient orthonormal basis, we may, when $F^\perp\neq0$, choose
$$
  e_0=x,
  \qquad
  e_1=\frac{F^\perp}{|F^\perp|},
  \qquad
  e_2,\ldots,e_{n-1}\perp\{x,F^\perp\}.
$$
These vectors are kept fixed as ambient vectors when differentiating the
component functions $F_r=\inner{F}{e_r}$. Put
$$
  z_r:=\nabla^S F_r(x)\in x^\perp,
  \qquad
  \beta_r=x\wedge z_r.
$$

Consider the Jacobi form on $x^\perp$,
$$
  \mathsf J_x^{\cE}(a,b)
  :=\cE(x\wedge a,x\wedge b).
$$
For $a\in x^\perp$ the sectional-curvature assumption gives
$$
  0\leq\mathsf J_x^{\cE}(a,a)\leq |a|^2,
$$
and hence
$$
  0\preceq\mathsf J_x^{\cE}\preceq\Id_{x^\perp}.
$$
Consequently, for all $a,z\in x^\perp$,
\begin{align}
  \mathsf J_x^{\cE}(z,z)-\mathsf J_x^{\cE}(a,z)
  &=
  \mathsf J_x^{\cE}(z-a/2,z-a/2)
  -\frac14\mathsf J_x^{\cE}(a,a) \geq-\frac14|a|^2.
  \label{eq:jacobi-complete-square}
\end{align}

We shall use the orthogonal-simple-form consequence
\eqref{eq:BK-simple-form} of Lemma~\ref{lem:BBK}, with $a=0$ and $b=1$.

We now estimate the summands
$$
  Q_r:=\cE(\beta_r,\beta_r)-\cE(\alpha_r,\beta_r).
$$
For $r=0$,
$$
  \alpha_0=F\wedge x=-x\wedge F^\perp,
$$
and therefore
$$
  Q_0
  =\mathsf J_x^{\cE}(z_0,z_0)
   -\mathsf J_x^{\cE}(-F^\perp,z_0)
  \geq-\frac14|F^\perp|^2
$$
by \eqref{eq:jacobi-complete-square}. Similarly,
$$
  \alpha_1=F\wedge e_1=\phi\,x\wedge e_1,
$$
because $F^\perp$ is parallel to $e_1$, and hence
$$
  Q_1
  =\mathsf J_x^{\cE}(z_1,z_1)
   -\mathsf J_x^{\cE}(\phi e_1,z_1)
  \geq-\frac14\phi^2.
$$

For $r\geq2$ we have
$$
  \alpha_r=\phi\,x\wedge e_r+F^\perp\wedge e_r.
$$
Thus
\begin{align*}
  Q_r
  ={}&
  \mathsf J_x^{\cE}(z_r,z_r)
  -\mathsf J_x^{\cE}(\phi e_r,z_r)
  -\cE(F^\perp\wedge e_r,\beta_r).
\end{align*}
The first two terms are bounded below by $-\phi^2/4$. Moreover,
$$
  \inner{F^\perp\wedge e_r}{x\wedge z_r}
  =
  \inner{F^\perp}{x}\inner{e_r}{z_r}
  -\inner{F^\perp}{z_r}\inner{e_r}{x}
  =0.
$$
Hence the two simple forms $F^\perp\wedge e_r$ and $\beta_r$ are
orthogonal, and \eqref{eq:BK-simple-form} gives
$$
  |\cE(F^\perp\wedge e_r,\beta_r)|
  \leq\frac23|F^\perp\wedge e_r|\,|\beta_r|
  =\frac23|F^\perp|\,|\beta_r|.
$$
It follows that
$$
  Q_r\geq-\frac14\phi^2
  -\frac23|F^\perp|\,|\beta_r|,
  \qquad r=2,\ldots,n-1.
$$

Summing the estimates for $r=0,\ldots,n-1$ yields the pointwise bound
\begin{equation}\label{eq:tangent-twist-pointwise}
  \sum_{r=0}^{n-1}
  \bigl\{\cE(\beta_r,\beta_r)-\cE(\alpha_r,\beta_r)\bigr\}
  \geq
  -\frac14\bigl(|F^\perp|^2+(n-1)\phi^2\bigr)
  -\frac23|F^\perp|
     \sum_{r=2}^{n-1}|\beta_r|.
\end{equation}
The same inequality holds when $F^\perp=0$ by continuity.

Multiplying by $c_{n,k}$ and integrating over $S(V)$, the first two
terms give
$$
  -\frac14\bigl(N_{\perp}+(n-1)N_{\parallel}\bigr).
$$
For the remaining term, Cauchy--Schwarz in the variables $(x,r)$ gives
\begin{align*}
  c_{n,k}\int_{S(V)}
  |F^\perp|\sum_{r=2}^{n-1}|\beta_r|\,d\sigma
  &\leq
  \left(
    c_{n,k}\int_{S(V)}
    \sum_{r=2}^{n-1}|F^\perp|^2\,d\sigma
  \right)^{1/2}\cdot
  \left(
    c_{n,k}\int_{S(V)}
    \sum_{r=2}^{n-1}|\beta_r|^2\,d\sigma
  \right)^{1/2}                                      \\
  &\leq
  \sqrt{(n-2)N_{\perp}}\,
  \left(
    c_{n,k}\int_{S(V)}
    \sum_{r=0}^{n-1}|\beta_r|^2\,d\sigma
  \right)^{1/2}.
\end{align*}
Although the adapted basis was chosen pointwise, the full sum
$\sum_r|\beta_r|^2$ is basis-independent. Using any fixed ambient
orthonormal basis,
$$
  \sum_{r=0}^{n-1}|\beta_r|^2
  =\sum_{r=0}^{n-1}|\nabla^S F_r|^2
  =|\nabla^S F|^2.
$$
Each $F_r$ is a spherical harmonic of degree $k$, since the corresponding
component of the homogeneous polynomial $F$ is harmonic. Therefore, with
$$
  \lambda_k=k(k+n-2),
$$
integration by parts on $S(V)$ gives
$$
  c_{n,k}\int_{S(V)}\sum_r|\beta_r|^2\,d\sigma
  =
  c_{n,k}\int_{S(V)}|\nabla^S F|^2\,d\sigma
  =
  \lambda_k c_{n,k}\int_{S(V)}|F|^2\,d\sigma
  =
  k(k+n-2)N.
$$
Substitution into \eqref{eq:tangent-twist-pointwise} proves
$$
  W_{\cE}(K)
  \geq-\frac14\bigl(N_{\perp}+(n-1)N_{\parallel}\bigr)
  -\frac23
   \sqrt{(n-2)N_{\perp}\,k(k+n-2)N},
$$
as claimed.
\end{proof}

We have the following positivity of the tangent-twist Pestov form.
\begin{lemma}
\label{lem:tangent-twist-high-degree}
Let $n\geq4$ and $k\geq5$. Suppose that the covariant tensor obtained from $K$ by identifying its
$V$-valued factor with $V^*$ via the metric, still denoted by $K$,
satisfies
\begin{equation}\label{eq:highest-normal-condition}
  \tf(\Sym K)=0.
\end{equation}
If
$
  \frac14\leq\sec_{\cA}\leq1,
$
then
$
  W_{\cA}(K)>0
$
for every $K\neq0$.
\end{lemma}

\begin{proof}
Assume that $K\neq0$. Let $\Harm_j(V)$ denote the space of homogeneous
harmonic polynomials of degree $j$ on $V$. Since the components of the
homogeneous extension of $F$ are harmonic of degree $k$,
$$
  \Delta_V\inner{F(x)}x
  =2\,\operatorname{div}_V F\in\Harm_{k-1}(V),
  \qquad
  \Delta_V^2\inner{F(x)}x=0.
$$
The Fischer decomposition of the homogeneous polynomial
$\inner{F(x)}x$ therefore has the form
$$
  \inner{F(x)}x=h_{k+1}(x)+|x|^2q(x),
  \qquad
  h_{k+1}\in\Harm_{k+1}(V),
  \quad
  q\in\Harm_{k-1}(V).
$$
The component $h_{k+1}$ is represented by the trace-free complete
symmetrization of $K$. Hence \eqref{eq:highest-normal-condition} gives
$$
  \inner{F(x)}x=|x|^2q(x),
  \qquad
  q\in\Sym^{k-1}_0V^*.
$$

The harmonic vector-valued polynomial with this radial part is
\begin{equation*}
  \widetilde F_q(x)
  =\frac{n+2k-4}{n+k-3}q(x)x
  -\frac{|x|^2}{n+k-3}\nabla^Vq(x).
\end{equation*}
Indeed, homogeneity of $q$ and the identity $\Delta_Vq=0$ give
$$
  \Delta_V\widetilde F_q=0,
  \qquad
  \inner{\widetilde F_q(x)}x=|x|^2q(x).
$$
Restricting to the unit sphere and using
$$
  \nabla^Vq=(k-1)qx+\nabla^Sq
$$
yields
$$
  F_q(x)
  =q(x)x-\frac1{n+k-3}\nabla^Sq(x).
$$

The contraction image generated by $q$ and the hook summand in the kernel
of radial contraction are orthogonal $O(n)$-types. Thus
\begin{equation}\label{eq:radial-contraction-decomposition}
  F=H+F_q,
  \qquad
  \inner{H(x)}x=0.
\end{equation}
Set
$$
  h:=c_{n,k}\int_{S(V)}|H|^2,
  \qquad
  s:=c_{n,k}\int_{S(V)}|F_q|^2.
$$
Since $q$ is a spherical harmonic of degree $k-1$,
$$
  \int_{S(V)}|\nabla^Sq|^2
  =(k-1)(n+k-3)\int_{S(V)}q^2.
$$
Using the orthogonality of the radial and tangential terms in $F_q$,
we obtain
\begin{equation}\label{eq:NPY-decomposition}
  N=h+s,
  \qquad
  N_{\parallel}
  =\frac{n+k-3}{n+2k-4}s,
  \qquad
  N_{\perp}
  =h+\frac{k-1}{n+2k-4}s.
\end{equation}
In particular, for
$
  u:={N_{\parallel}} / {N},
$
we have
\begin{equation*}
  0\leq u\leq\frac{n+k-3}{n+2k-4}<1,
  \qquad
  \frac{N_{\perp}}N=1-u.
\end{equation*}

Put
$$
  \lambda:=k(k+n-2).
$$
For the constant-curvature operator $I$, direct calculation on the two
summands in \eqref{eq:radial-contraction-decomposition} gives
\begin{equation*}
  W_I(K)
  =(\lambda-1)h+(k-1)(k+n-2)s.
\end{equation*}
Equivalently, using \eqref{eq:NPY-decomposition},
\begin{equation}\label{eq:WI-u}
  \frac{W_I(K)}N
  =\lambda-1-(n+2k-4)u.
\end{equation}
Indeed, the Jacobi contribution is $\lambda N$, whereas the twist
contribution is
$$
  -N-(n+2k-4)N_{\parallel}.
$$

Write
$$
  \cA=\frac14I+\frac34\cE,
  \qquad
  0\leq\sec_{\cE}\leq1.
$$
By linearity,
$$
  W_{\cA}(K)=\frac14W_I(K)+\frac34W_{\cE}(K).
$$
Combining \eqref{eq:WI-u} with
Lemma \ref{lem:tangent-twist-residual} and using
$N_{\perp}/N=1-u$, we obtain the stronger estimate
\begin{equation}\label{eq:G-lower-bound}
  \frac{W_{\cA}(K)}N\geq G_{n,k}(u),
  \qquad
  0\leq u\leq\frac{n+k-3}{n+2k-4},
\end{equation}
where
\begin{equation}\label{eq:Gnk}
  G_{n,k}(u)
  :=\frac{4\lambda-7}{16}
  -\frac{7n+8k-22}{16}u
  -\frac12\sqrt{(n-2)\lambda(1-u)}.
\end{equation}

It remains to verify that this lower bound is positive. Set
$$
  L:=7n+8k-22,
  \qquad
  Q:=(n-2)\lambda,
  \qquad
  z:=\sqrt{1-u}.
$$
Completing the square gives
\begin{equation}\label{eq:G-completed-square}
  G_{n,k}(u)
  =\frac{\Pi_{n,k}}{16L}
   +\frac{L}{16}
    \left(z-\frac{4\sqrt Q}{L}\right)^2,
\end{equation}
where
\begin{align}
  \Pi_{n,k}
  ={}&(12k-49)n^2
     +(44k^2-192k+259)n \notag\\
    &+32k^3-184k^2+408k-330.
  \notag
\end{align}
For $k\geq5$, all three coefficients of this polynomial in $n$ are
positive. Indeed,
$$
  12k-49\geq11,
$$
the second coefficient is increasing for $k\geq5$ and equals $399$ at
$k=5$, while the third is increasing and equals $1110$ at $k=5$.
Thus $\Pi_{n,k}>0$. Since $L>0$, \eqref{eq:G-completed-square} implies
$$
  G_{n,k}(u)>0.
$$
The conclusion now follows from \eqref{eq:G-lower-bound}.
\end{proof}

\subsubsection{Reduction to degrees one and three}

\begin{proposition}[Degree reduction]\label{prop:degree-reduction}
Let $n\geq4$ be even and assume that $g$ is strictly $1/4$-pinched. Then any
odd invariant unit normal vector field
$f\in C^\infty(\SM,\cN)$ with $Xf=0$ has the form
$$
  f=f_1+f_3.
$$
If $f_3=0$, then no such field exists.
\end{proposition}

\begin{proof}
Choose a pinching presentation with $\delta>1/4$ and let $\cA$ be the
normalized curvature operator from \eqref{eq:normalized-curvature}. Let $k$
be the highest Fourier degree of $f$, and denote its top coefficient by
$K_k\in\Sym^k_0T^*M\otimes TM$. Since $f$ is odd, $k$ is odd, and
$K_k\neq0$. The top-degree part of $Xf=0$ is $X_+f_k=0$, so $K_k$ is a
twisted conformal Killing tensor. Moreover, the degree-$(k+1)$ harmonic
component of the identity
$$
  \inner{f(x,v)}{v}=0
$$
comes only from $f_k$ and, under the harmonic--tensor correspondence, is
precisely $\tf(\Sym K_k)$. Thus $K_k$ satisfies
\eqref{eq:highest-normal-condition}.

If $k\geq5$, then $\frac14\leq\sec_{\cA}\leq1$, and
Lemma \ref{lem:tangent-twist-high-degree} gives
$$
  W_{\cA}(K_k)>0
$$
at every point where $K_k\neq0$. Since $K_k$ is a nonzero smooth section,
its integral is therefore strictly positive, contradicting
\eqref{eq:tangent-twist-pestov-normalized}. Hence $k\leq3$, and oddness gives
$f=f_1+f_3$.

Assume now that $f_3=0$. Then $f(x,v)=J_xv$ for an endomorphism field
$J\in C^\infty(M,\End(TM))$. The normality and unit-length conditions imply
$$
  \inner{Jv}{v}=0,
  \qquad
  |Jv|=|v|
$$
for every unit vector $v$, and hence, by homogeneity, for every $v\in TM$.
Polarizing the first identity gives
$$
  \inner{Ju}{w}+\inner{u}{Jw}=0,
$$
while polarizing the second gives
$$
  \inner{Ju}{Jw}=\inner{u}{w}.
$$
Thus $J^*=-J$ and $J^*J=\Id$, so in particular $J^2=-\Id$. Therefore $J$
is an orthogonal almost complex structure.

For a geodesic $\gamma$ with $\dot\gamma(0)=v$, the equation $Xf=0$ gives
$$
  0=\left.\nabla_{\dot\gamma}
     \bigl(J_{\gamma(t)}\dot\gamma(t)\bigr)\right|_{t=0}
   =(\nabla_vJ)v,
$$
since $\dot\gamma$ is parallel. Hence $(M,g,J)$ is nearly K\"ahler.
Nagy's decomposition theorem \cite{Nagy02} splits the universal cover into
a K\"ahler factor and a strict nearly K\"ahler factor. Negative sectional
curvature excludes a non-trivial Riemannian product, since mixed two-planes
in such a product have zero sectional curvature. Thus the universal cover
is either K\"ahler or strict nearly K\"ahler. The latter has positive scalar
curvature \cite{Gray76,Nagy02}, contradicting negative sectional curvature,
whereas a negatively curved K\"ahler metric cannot be strictly more than
quarter-pinched \cite{Berger60}; see also
\cite[proof of Theorem~4.1]{CLMS24}.
\end{proof}

\subsection{The \texorpdfstring{$(3,1)$}{(3,1)}-type tensor and Weitzenb\"ock identities}\label{sec:cubic}

\subsubsection{From the invariant field to a generalized Killing tensor of Young type \texorpdfstring{$(3,1)$}{(3,1)}}

From now on, let
$$
  f=f_1+f_3
$$
be an odd invariant normal field. Under the standard harmonic--tensor
identification there are unique fields
$$
  A\in C^\infty(M,T^*M\otimes TM),
  \qquad
  K\in C^\infty(M,\Sym^3_0T^*M\otimes TM)
$$
such that
$$
  f_1(x,v)=A_xv,
  \qquad
  f_3(x,v)=K_x(v,v,v).
$$
Using $g$, we identify the $TM$-valued factors of $A$ and $K$ with
$T^*M$, and use the same letters for the vector-valued tensors and
their covariant counterparts.

\begin{definition}\label{def:T}
Define $T\in C^\infty(M,\Sym^3T^*M\otimes T^*M)$ by
\begin{align}
  T(u_1,u_2,u_3,w)
  :={}&K(u_1,u_2,u_3,w)\notag\\
  &+\frac13\bigl(
      g(u_1,u_2)A(u_3,w)
     +g(u_1,u_3)A(u_2,w)
     +g(u_2,u_3)A(u_1,w)
    \bigr).
  \label{eq:Tdef}
\end{align}
\end{definition}
Then
\begin{equation}\label{eq:diagonal-T}
  T(v,v,v,w)=\inner{f(x,v)}{w}
  \qquad\text{for }|v|=1,
\end{equation}
and, by homogeneity,
$$
  T(v,v,v,w)
  =\inner{|v|^2Av+K(v,v,v)}{w}
$$
for arbitrary $v$.


\begin{lemma}[Young symmetry]\label{lem:young-symmetry}
The normality condition $\inner{f(x,v)}{v}=0$ is equivalent to
$
  T_{(abcd)}=0.
$
Moreover,
$$
  T_x\in\cC_x
  :=\ker\bigl(
    \Sym^3T_x^*M\otimes T_x^*M
    \xrightarrow{\Sym}
    \Sym^4T_x^*M
  \bigr).
$$
As a $GL(T_xM)$-module,
$$
  \cC_x\simeq\Hook T_x^*M.
$$
\end{lemma}

\begin{proof}
By \eqref{eq:diagonal-T}, normality gives
$$
  T(v,v,v,v)=0
$$
for all unit vectors. Homogeneity gives the same identity for every vector.
A symmetric four-linear form is determined by its diagonal polynomial, so the
complete symmetrisation of $T$ vanishes. Conversely, vanishing of the complete
symmetrisation gives $T(v,v,v,v)=0$ and hence normality.

The Pieri decomposition
$$
  \Sym^3V^*\otimes V^*
  \simeq \Sym^4V^*\oplus\Hook V^*
$$
shows that the kernel is precisely the Schur module of shape $(3,1)$; see
\cite[Chapter~6]{FultonHarris}.
\end{proof}

We henceforth identify $\cC$ with the hook bundle $\Hook T^*M$.

\begin{lemma}[Generalized Killing equation]\label{lem:killing-equation}
The flow-invariance equation $Xf=0$ is equivalent to
\begin{equation}\label{eq:generalized-killing}
  \nabla_{(e}T_{abc)d}=0.
\end{equation}
\end{lemma}

\begin{proof}
Let $\gamma$ be a geodesic with $\dot\gamma(0)=v$, and let $w(t)$ be parallel
along $\gamma$. Since $\dot\gamma$ is parallel,
$$
  \frac{d}{dt}
  T_{\gamma(t)}(\dot\gamma,\dot\gamma,\dot\gamma,w(t))
  =(\nabla_vT)(v,v,v,w)
\quad\text{at }t=0.
$$
Hence $Xf=0$ implies
$$
  (\nabla_vT)(v,v,v,w)=0
$$
for every $v,w$. Polarising the degree-$4$ polynomial in $v$ gives
\eqref{eq:generalized-killing}. The converse follows by restriction to the diagonal.
\end{proof}

Equivalently, if
$$
  \cD:C^\infty(M,\cC)
  \longrightarrow C^\infty(M,\Sym^4T^*M\otimes T^*M)
$$
is defined by
$
  (\cD T)_{eabcd}:=\nabla_{(e}T_{abc)d},
$
then $\cD T=0$.
\subsubsection{A second Weitzenb\"ock identity on the totally trace-free \texorpdfstring{$(3,1)$}{(3,1)}-component}\label{sec:tracefree-hook-weitzenbock}

The ordinary formula of Appendix \ref{sec:weitzenbock} is sufficient in dimension
$4$, but its curvature term is not pointwise positive on the full
$(3,1)$-module in higher dimension.  The purpose of this subsection is to
extract a second formula on the totally trace-free $(3,1)$-component.  Throughout this
subsection $V$ is a Euclidean vector space of dimension $n\geq8$, and
$$
  \cU:=\mathbb V^0_{(3,1)}V^*,
$$
where $\mathbb V^0_\lambda V^*$ denotes the irreducible trace-free
$O(V)$-module of highest weight $\lambda$. Restriction from $GL(V)$ to
$O(V)$ gives the parallel orthogonal
decomposition
\begin{equation}\label{eq:orthogonal-hook-decomposition}
  \Hook V^*
  =\cU\oplus\Sym^2_0V^*\oplus\Lambda^2V^*.
\end{equation}
Write $\operatorname{pr}_{\cU}$ for the first projection. If
$X\in\mathfrak{so}(V)$, decompose the infinitesimal action on a hook as
$$
  \rho_X=\mathsf D_X+\mathsf O_X,
$$
where $\mathsf D_X$ is the sum of the actions in the first three symmetric
slots and $\mathsf O_X$ is the action in the fourth slot.

\paragraph{The cubic formula.}

In the $\Lambda^2V$-normalization of
\cite[\S3.2]{SemmelmannWeingart10}, the Casimir eigenvalue of
$\cU$ is $-4n$. Substitution into \cite[(4.18)]{SemmelmannWeingart10}
gives the odd cubic Weitzenb\"ock polynomial
\begin{equation}\label{eq:p3-polynomial}
  p_3(z)=z^3+(n-1)z^2+
  \left(\frac{n(n-2)}4-8\right)z-4n.
\end{equation}
Put
\begin{equation}\label{eq:partial-formula}
  \beta_n:=\frac{n^2-10n-8}{4},
  \qquad
  F_{\cU}:=\frac{p_3(B)-\beta_nB}{4n}
  \quad\text{on }V^*\otimes\cU.
\end{equation}

For an algebraic curvature operator $\cB$ on $V$, define the quadratic form
\begin{equation}\label{eq:wR-definition}
  w_{\cB}(h,h)
  :=\sum_{\alpha,\beta}\cB_{\alpha\beta}
  \inner{\mathsf D_{E_\alpha}h}{(\mathsf D_{E_\beta}+\mathsf O_{E_\beta})h},
  \qquad
  \cB_{\alpha\beta}:=\inner{\cB E_\alpha}{E_\beta}.
\end{equation}
The definition is independent of the orthonormal basis $(E_\alpha)$.

\begin{lemma}\label{lem:first-three-slot-action}
For $a\in V^*$ and $h\in\cU$,
\begin{equation}\label{eq:first-three-slot-action}
  F_{\cU}(a\otimes h)
  =\sum_i e^i\otimes\operatorname{pr}_{\cU}
     \mathsf D_{e_i\wedge a}h.
\end{equation}
\end{lemma}

\begin{proof}
We give the contraction which fixes the normalisation. For
$u\in\Sym^3V^*\otimes V^*$, let $\Sym_4$ denote normalized symmetrization
over its $4$ tensor slots and set
$$
  \mathsf Yu:=u-\Sym_4u,
  \qquad C(\mathsf Yu)_{cd}:=(\mathsf Yu)_{aa,cd},
$$
and define
$$
  C^*q:=\mathsf Y\!\left[\frac13
  (g_{ab}q_{cd}+g_{ac}q_{bd}+g_{bc}q_{ad})\right].
$$
For a two-tensor $q$, write
$q_{\mathrm s}:=(q+q^\top)/2$ and $q_{\mathrm a}:=(q-q^\top)/2$.
A direct trace gives
$$
  CC^*|_{\Sym^2_0}=\frac n6\Id,
  \qquad
  CC^*|_{\Lambda^2}=\frac{n+2}{3}\Id,
$$
and hence
\begin{equation}\label{eq:tracefree-hook-projection}
  \operatorname{pr}_{\cU}u
  =\mathsf Yu-\frac6nC^*((C\mathsf Yu)_{\mathrm s})
     -\frac3{n+2}C^*((C\mathsf Yu)_{\mathrm a}).
\end{equation}
Write $\rho_{ab}:=\rho(e_a\wedge e_b)$ and
$$
  \mathsf M_{ab}:=\operatorname{pr}_{\cU}\mathsf D_{e_a\wedge e_b}
                  \operatorname{pr}_{\cU}.
$$
For $m\geq1$, let
$$
  (B^m)_{ab}h
  :=(e_a\mathbin{\lrcorner}\otimes\Id)B^m(e^b\otimes h),
$$
and write $\operatorname{Alt}_{ab}Q_{ab}:=(Q_{ab}-Q_{ba})/2$.
Expanding $B^2$ and $B^3$ by \eqref{eq:Bformula}, applying
\eqref{eq:tracefree-hook-projection} after each word in $\mathsf D$ and
$\mathsf O$, and using the Young relation and the vanishing traces gives
\begin{align}
 \operatorname{Alt}_{ab}(B^2)_{ab}
   &=-\frac{n-2}{2}\rho_{ab},\label{eq:alt-B2}\\
 \operatorname{Alt}_{ab}(B^3)_{ab}
   &=4n\mathsf M_{ab}+\frac{n^2-7n+14}{2}\rho_{ab}.
   \label{eq:alt-B3}
\end{align}
The word-by-word coefficients are recorded in
Appendix~\ref{app:first-three-slot}. Since $p_3$ is twist-odd, only these alternating
parts occur. The coefficient of $\rho_{ab}$ in $p_3(B)$ is
$$
 \frac{n^2-7n+14}{2}
 -\frac{(n-1)(n-2)}2+\frac{n(n-2)}4-8
 =\frac{n^2-10n-8}{4}=\beta_n.
$$
Consequently, for $a\otimes h\in V^*\otimes\cU$,
$$
  p_3(B)(a\otimes h)
  =4n\sum_i e^i\otimes\operatorname{pr}_{\cU}
       \mathsf D_{e_i\wedge a}h+\beta_nB(a\otimes h),
$$
which proves \eqref{eq:first-three-slot-action}.
\end{proof}

Orthogonal Pieri gives
\begin{equation}\label{eq:Vtensor-tracefree-hook}
\begin{split}
 V^*\otimes\cU={}&
 \mathbb V^0_{(4,1)}\oplus\mathbb V^0_{(3,2)}
 \oplus\mathbb V^0_{(3,1,1)}\oplus\Sym^3_0V^*\oplus\mathbb V^0_{(2,1)}.
\end{split}
\end{equation}

\begin{lemma}\label{lem:first-three-slot-spectrum}
On the five summands in \eqref{eq:Vtensor-tracefree-hook}, in the displayed
order, the eigenvalues of $B$ and $F_{\cU}$ are
\begin{equation}\label{eq:partial-spectrum-table}
\begin{array}{c|ccccc}
 & (4,1)&(3,2)&(3,1,1)&(3)&(2,1)\\ \hline
 B&3&0&-2&2-n&-n-1\\[2pt]
 F_{\cU}&\dfrac{11}{4}&-1&-1&
 -\dfrac{n^2-2n+8}{4n}&-\dfrac{n^2+n-1}{n}.
\end{array}
\end{equation}
In particular $F_{\cU}\preceq-\Id$ off the $(4,1)$ summand.

If $T\in\Hook V^*$ satisfies the generalized Killing symbol
equation and $H=\operatorname{pr}_{\cU}T$, then the $(4,1)$ component of $\nabla H$ vanishes.
\end{lemma}

\begin{proof}
The conformal weights in the first row of
\eqref{eq:partial-spectrum-table} are the standard orthogonal Pieri weights.
Substitution into \eqref{eq:p3-polynomial}--\eqref{eq:partial-formula}
gives the second row.  The only positive eigenvalue is therefore the one on
$\mathbb V^0_{(4,1)}$.

The generalized Killing symbol is the $GL(V)$-projection onto
$\mathbb S_{(4,1)}V^*$.  The trace-free $O(V)$-type
$\mathbb V^0_{(4,1)}$ occurs in $V^*\otimes\cU$, but not in
$$
 V^*\otimes\Sym^2_0V^*
 =\Sym^3_0V^*\oplus\mathbb V^0_{(2,1)}\oplus V^*
$$
or in
$$
 V^*\otimes\Lambda^2V^*
 =\mathbb V^0_{(2,1)}\oplus\Lambda^3V^*\oplus V^*.
$$
It therefore cannot be cancelled by either trace block, and its projection
in $\nabla H$ is zero.\end{proof}

Let now $T$ be a smooth generalized Killing tensor of Young type $(3,1)$ on a closed manifold and
$H=\operatorname{pr}_{\cU}T$. Let $\widetilde F_{\cU}$ be the contraction associated with
$F_{\cU}$ as in Appendix \ref{sec:weitzenbock}. By
\eqref{eq:first-three-slot-action}, the skew-adjointness of the infinitesimal
orthogonal action gives the pointwise identity
\begin{equation*}
  \inner{\widetilde F_{\cU}(\nabla^2H)}{H}
  =-w_{\Rop^g}(H,H).
\end{equation*}
On the other hand, integration by parts gives
$$
  \int_M\inner{\widetilde F_{\cU}(\nabla^2H)}{H}\,d\vol_g
  =-\int_M\inner{F_{\cU}\nabla H}{\nabla H}\,d\vol_g.
$$
By Lemma \ref{lem:first-three-slot-spectrum}, the $(4,1)$ component of $\nabla H$ vanishes
and $F_{\cU}\preceq-\Id$ on all remaining components. Therefore
\begin{equation}\label{eq:partial-integral-sign}
  \int_M w_{\Rop^g}(H,H)\,d\vol_g
  =\int_M\inner{F_{\cU}\nabla H}{\nabla H}\,d\vol_g
  \leq0.
\end{equation}

\paragraph{The pointwise twisted estimate.}
We have the tangent-twist estimate on the totally trace-free $(3,1)$-component.
\begin{lemma}\label{lem:tracefree-hook-twist-estimate}
Let $n\geq8$ be even, let $H\in\mathbb V^0_{(3,1)}V^*$, and let $\cE$ be an
algebraic curvature tensor satisfying $0\leq\sec_{\cE}\leq1$.  Then
\begin{equation}\label{eq:tracefree-hook-twist-bound}
  w_{\cE}(H,H)\geq-B_n|H|^2,
  \qquad
  B_n:=\frac12\sqrt{\frac{n(n-2)(n+2)}{n-1}}
       +\frac23\sqrt{\frac{n(n-2)}2}.
\end{equation}
Moreover,
\begin{equation}\label{eq:wI-tracefree-hook}
  w_I(H,H)=(3n+2)|H|^2.
\end{equation}
Consequently, for $\cA=\frac14I+\frac34\cE$,
\begin{equation}\label{eq:tracefree-hook-twist-positive}
  w_{\cA}(H,H)\geq c_n|H|^2,
  \qquad
  c_n:=\frac{3n+2-3B_n}{4}>0.
\end{equation}
\end{lemma}

\begin{proof}
Use the spherical model
$$
  F(x):=H(x,x,x,\cdot)^\sharp,
  \qquad x\in S^{n-1},
$$
and choose $\kappa_n>0$ so that
$|H|^2=\kappa_n\int_{S^{n-1}}|F|^2$.  The totally trace-free $(3,1)$ conditions give
\begin{equation}\label{eq:pure-spherical-conditions}
  \inner{F(x)}x=0,
  \qquad \operatorname{div}_{S^{n-1}}F=0.
\end{equation}
At a fixed $x$, put $e_0=x$ and choose an orthonormal basis
$e_1,\ldots,e_{n-1}$ of $x^\perp$.  Let
$$
  \mathsf A(u):=\nabla^S_uF,
  \qquad \mathsf A=\Sigma+\Omega,
$$
be the symmetric/skew decomposition. Then $\tr\mathsf A=0$.  For the fixed
ambient components $F_r=\inner F{e_r}$, set
$$
  \alpha_r=F\wedge e_r,
  \qquad \beta_r=x\wedge\nabla^SF_r.
$$
Directly expanding \eqref{eq:wR-definition} gives
\begin{equation}\label{eq:spherical-w-formula}
 w_{\cE}(H,H)=\kappa_n\int_{S^{n-1}}
 \sum_{r=0}^{n-1}
 \{\cE(\beta_r,\beta_r)-\cE(\alpha_r,\beta_r)\}\,d\sigma(x).
\end{equation}
The $r=0$ term vanishes, because
$\nabla^SF_0=-F$ and hence $\alpha_0=\beta_0=F\wedge x$.  For $a\geq1$,
$$
 \beta_a=x\wedge\mathsf A^*e_a,
 \qquad \alpha_a=F\wedge e_a.
$$
The first part of every remaining summand in
\eqref{eq:spherical-w-formula} is nonnegative.  Thus it remains to bound
$$
 \mathfrak m:=\sum_{a=1}^{n-1}\cE(F\wedge e_a,x\wedge\mathsf A^*e_a).
$$
Diagonalise $\Sigma$, with eigenvalues $s_a$.  The first Bianchi identity gives
\begin{equation}\label{eq:M-decomposition}
 \mathfrak m=\sum_a s_a\cE(F,e_a,x,e_a)
   +\cE(F\wedge x,\omega_\Omega),
 \qquad
 \omega_\Omega:=\sum_{a<b}\Omega_{ab}e_a\wedge e_b.
\end{equation}
For each unit $e$, the Jacobi form
$\mathsf J_e^{\cE}(u,v):=\cE(u,e,v,e)$ satisfies
$0\preceq\mathsf J_e^{\cE}\preceq\Id$.  Since the
relevant vectors are orthogonal, centering at $\frac12\Id$ gives
$$
  |\cE(F,e_a,x,e_a)|\leq\frac12|F|.
$$
If $m=n-1$ and $\sum_as_a=0$, then, because $m$ is odd,
\begin{equation}\label{eq:tracefree-l1}
  \sum_a|s_a|\leq\chi_m|\Sigma|,
  \qquad \chi_m^2=\frac{n(n-2)}{n-1}.
\end{equation}
Indeed, if $p$ and $q$ are the numbers of positive and negative eigenvalues,
respectively, trace-freeness and Cauchy--Schwarz on the two sign classes give
$$
 \left(\sum_a|s_a|\right)^2
 \leq\frac{4pq}{p+q}|\Sigma|^2.
$$
Zero eigenvalues may be omitted.  Put $r=p+q\leq m=n-1$.  If $r=m$,
oddness gives $4pq\leq m^2-1$; if $r\leq m-1$, then
$4pq/r\leq r\leq m-1$.  In either case
$4pq/(p+q)\leq(m^2-1)/m=n(n-2)/(n-1)$, proving
\eqref{eq:tracefree-l1}.
Put $\omega_\Omega$ in canonical form.  Its rank is at most $n-2$, and
each of its simple summands is orthogonal to the simple form $F\wedge x$.
Applying Lemma~\ref{lem:BBK} to the orthogonal simple summands in a
canonical form of $\omega_\Omega$ gives
\begin{equation*}
 |\cE(F\wedge x,\omega_\Omega)|
 \leq\frac23\sqrt{\frac{n-2}{2}}\,|F|\,|\omega_\Omega|.
\end{equation*}
It remains to insert the exact spherical energies.  Each fixed ambient
component of $F$ is a degree-three spherical harmonic, whose eigenvalue is
$\lambda_3=3(n+1)$.  Since the ambient derivative of the tangent field is
$$
 \nabla^V_uF=\nabla^S_uF-\inner{F}{u}x,
$$
we obtain
$$
 \kappa_n\int|\mathsf A|^2
 =\{3(n+1)-1\}|H|^2=(3n+2)|H|^2.
$$
Let $\eta=F^\flat$. By \eqref{eq:pure-spherical-conditions},
$d^*\eta=0$, while $d\eta=2\omega_\Omega$ up to the immaterial sign fixed
by the matrix convention.  The Hodge--Bochner identity on $S^{n-1}$,
whose Ricci tensor is $(n-2)g$, therefore gives
$$
 4\kappa_n\int|\omega_\Omega|^2
 =\kappa_n\int|d\eta|^2
 =\kappa_n\int\bigl(|\mathsf A|^2+(n-2)|F|^2\bigr)
 =4n|H|^2.
$$
Finally, our norm conventions are
$$
 |\Omega|_{\End}^2=2|\omega_\Omega|^2,
 \qquad |\mathsf A|^2=|\Sigma|^2+2|\omega_\Omega|^2.
$$
Thus
\begin{equation}\label{eq:tracefree-hook-energies}
\begin{split}
 \kappa_n\int|\Sigma|^2&=(n+2)|H|^2,\\
 \kappa_n\int|\omega_\Omega|^2&=n|H|^2,\\
 \kappa_n\int|\mathsf A|^2&=(3n+2)|H|^2.
\end{split}
\end{equation}
Combining \eqref{eq:M-decomposition}--\eqref{eq:tracefree-hook-energies} and
Cauchy--Schwarz yields exactly
$$
 \kappa_n\int|\mathfrak m|
 \leq\left\{
 \frac12\sqrt{\frac{n(n-2)(n+2)}{n-1}}
 +\frac23\sqrt{\frac{n(n-2)}2}\right\}|H|^2.
$$
This proves \eqref{eq:tracefree-hook-twist-bound}.  The last identity in
\eqref{eq:tracefree-hook-energies}, or a direct Casimir calculation, gives
\eqref{eq:wI-tracefree-hook}.

Finally,
$$
 \frac12\sqrt{\frac{n(n-2)(n+2)}{n-1}}<\frac n2+\frac14,
 \qquad
 \frac23\sqrt{\frac{n(n-2)}2}<\frac{\sqrt2}{3}(n-1),
$$
and hence
$$
 n+\frac23-B_n
 >n\left(\frac12-\frac{\sqrt2}{3}\right)
   +\frac5{12}+\frac{\sqrt2}{3}>0.
$$
This proves \eqref{eq:tracefree-hook-twist-positive}.
\end{proof}


\subsection{Dimensions 4, 6, and the stable range}\label{sec:dimension-estimates}

The stable-range argument is contained in Proposition \ref{prop:stable-cubic}. We now
prove the two low-dimensional estimates needed for dimensions $4$ and $6$.

\subsubsection{The 4-dimensional curvature estimate}\label{sec:curvature}

Let $V$ be an oriented Euclidean $4$-space. Denote by
$$
  *:\Lambda^2V\longrightarrow\Lambda^2V
$$
the Hodge star. We prove the pointwise inequality needed to contradict
\eqref{eq:integral-negative-Q}.

\paragraph{A quantitative Finsler--Thorpe lemma.}

\begin{lemma}[Finsler--Thorpe with a parameter bound]\label{lem:FT}
Let $\cA:\Lambda^2V\to\Lambda^2V$ be an algebraic curvature operator
satisfying
$$
  \delta\leq\sec_{\cA}\leq1
$$
for some $\delta>1/4$. Then there exists $t\in(-1/2,1/2)$ such that
\begin{equation}\label{eq:FTpositive}
  \cP:=\cA-\frac14I+t*\succeq0.
\end{equation}
\end{lemma}

\begin{proof}
Set $\cS:=\cA-\frac14I$. Then
$$
  \delta-\frac14\leq\sec_{\cS}\leq\frac34,
$$
so $\cS$ has positive sectional curvature. The Finsler--Thorpe lemma states that
$$
  \mathcal I(\cS):=\{t\in\bbR:\cS+t*\succeq0\}
$$
is a nonempty closed interval; see, for example,
\cite[Section~2]{BettiolGoodman24}. Write $\mathcal I(\cS)=[t_-,t_+]$.

At the lower endpoint there is a nonzero
$\alpha\in\ker(\cS+t_-*)$ with $\inner{*\alpha}{\alpha}>0$. To justify the sign,
note that increasing $t$ moves into the positive-semidefinite region. If the
star quadratic form vanished on a kernel vector, that vector would be a
simple two-form in dimension $4$; then
$\inner{\cS\alpha}{\alpha}=0$, contradicting the strict positivity of sectional
curvature. After an oriented orthogonal change of basis,
$$
  \alpha=a e_{12}+b e_{34},
  \qquad a,b>0.
$$
Put
$$
  x=\cS_{1212},\qquad y=\cS_{3434},\qquad z=\cS_{1234}.
$$
The equations obtained by projecting $(\cS+t_-*)\alpha=0$ onto
$e_{12}$ and $e_{34}$ give
$$
  xa+(z+t_-)b=0,
  \qquad
  (z+t_-)a+yb=0.
$$
Hence $z+t_-=-\sqrt{xy}$ and therefore
$$
  t_-=-z-\sqrt{xy}\leq|z|.
$$
At the upper endpoint the analogous kernel vector has the form
$a e_{12}-b e_{34}$, and one obtains
$$
  t_+=-z+\sqrt{xy}\geq-|z|.
$$

Lemma~\ref{lem:BBK}, applied to the orthogonal simple forms
$e_{12}$ and $e_{34}$ and to the sectional spread of $\cS$, gives
$$
  |z|
  \leq\frac23\bigl(\max\sec_{\cS}-\min\sec_{\cS}\bigr)
  =\frac23(1-\delta)<\frac12.
$$
Here the extrema are taken over unit simple two-forms.
Thus $t_-<1/2$ and $t_+>-1/2$. Since $t_-\leq t_+$, the interval
$\mathcal I(\cS)$ meets $(-1/2,1/2)$. Any $t$ in the intersection satisfies
\eqref{eq:FTpositive}.
\end{proof}

\paragraph{The \texorpdfstring{$\Spin(4)$}{Spin(4)} decomposition.}

Use the double cover
$$
  \Spin(4)=\SU(2)_+\times\SU(2)_-.
$$
Let
$$
  V_{p,q}:=\Sym^p\bbC^2\boxtimes\Sym^q\bbC^2.
$$
The complexified standard representation is $V_{1,1}$. Since
$$
  \Sym^3V_{1,1}=V_{3,3}\oplus V_{1,1}
$$
and
$$
  \Sym^4V_{1,1}=V_{4,4}\oplus V_{2,2}\oplus V_{0,0},
$$
the identity
$$
  \Sym^3V_\bbC^*\otimes V_\bbC^*
  =\Sym^4V_\bbC^*\oplus\Hook V_\bbC^*
$$
gives
\begin{equation}\label{eq:spin-decomposition}
  E_\bbC
  \simeq
  V_{4,2}\oplus V_{2,4}\oplus V_{2,2}
  \oplus V_{2,0}\oplus V_{0,2}.
\end{equation}
The last two summands are the complexifications of
$\Lambda^2_+V$ and $\Lambda^2_-V$.

Let $C_+$ and $C_-$ be the positive Casimir endomorphisms of the two
$\mathfrak{su}(2)$ factors. Up to the same positive scalar normalisation,
\begin{equation*}
  \cQ_E(I)=C_++C_-,
  \qquad
  \cQ_E(*)=C_+-C_-.
\end{equation*}
On $V_{p,q}$, denote the scalar ratio by
\begin{equation}\label{eq:casimir-ratio}
  \eta_{p,q}:=\frac{\cQ_E(*)}{\cQ_E(I)}
  =\frac{p(p+2)-q(q+2)}{p(p+2)+q(q+2)}.
\end{equation}
In particular,
\begin{equation*}
  \eta_{4,2}=\frac12,
  \qquad
  \eta_{2,4}=-\frac12,
  \qquad
  \eta_{2,2}=0,
  \qquad
  \eta_{2,0}=1,
  \qquad
  \eta_{0,2}=-1.
\end{equation*}

\paragraph{Positivity on the first three summands.}

Let $t$ and $\cP$ be supplied by Lemma \ref{lem:FT}. Since
$$
  \cA=\cP+\frac14I-t*,
$$
linearity gives
\begin{equation}\label{eq:Qsplit}
  \cQ_E(\cA)
  =\cQ_E(\cP)+\frac14\cQ_E(I)-t\cQ_E(*).
\end{equation}
The first term is nonnegative by \eqref{eq:Qpositive}. On $V_{4,2}$,
\eqref{eq:casimir-ratio} gives
$$
  \cQ_E(\cA)
  \succeq
  \left(\frac14-\frac t2\right)\cQ_E(I)
  \succ0,
$$
because $t<1/2$. Similarly,
$$
  \cQ_E(\cA)|_{V_{2,4}}
  \succeq
  \left(\frac14+\frac t2\right)\cQ_E(I)
  \succ0,
$$
and
$$
  \cQ_E(\cA)|_{V_{2,2}}
  \succeq\frac14\cQ_E(I)
  \succ0.
$$

\paragraph{The two-form summands.}

It remains to treat $V_{2,0}\oplus V_{0,2}$, which is the ordinary two-form
representation. Every real two-form can be written in an oriented orthonormal
basis as
$$
  \omega=a e_{12}+b e_{34}.
$$
A direct calculation of the curvature term gives
\begin{equation}\label{eq:two-form-Q}
  \inner{\cQ_{\Lambda^2}(\cA)\omega}{\omega}
  =\Xi_{\cA}(a^2+b^2)-4ab\,\cA_{1234},
\end{equation}
where
$$
  \Xi_{\cA}:=\cA_{1313}+\cA_{1414}+\cA_{2323}+\cA_{2424}.
$$
Indeed, this is the usual Bochner curvature term on two-forms; for constant
curvature one it reduces to $4|\omega|^2$, as it should in dimension $4$.

By the sectional lower bound, $\Xi_{\cA}\geq4\delta$.  Lemma~\ref{lem:BBK}
gives
$$
  |\cA_{1234}|\leq\frac23(1-\delta).
$$
Using $4|ab|\leq2(a^2+b^2)$ in \eqref{eq:two-form-Q}, we obtain
\begin{align}
  \inner{\cQ_{\Lambda^2}(\cA)\omega}{\omega}
  &\geq
  \left(4\delta-2|\cA_{1234}|\right)|\omega|^2\notag\\
  &\geq
  \frac{4(4\delta-1)}3|\omega|^2>0.
  \label{eq:two-form-positive}
\end{align}

Combining the preceding subsections gives the central algebraic statement.

\begin{proposition}
\label{prop:curvature-positive}
Let $V$ be an oriented Euclidean $4$-space and let $\cA$ be an algebraic
curvature operator satisfying
$$
  \frac14<\sec_{\cA}\leq1.
$$
Then
$$
  \cQ_{\Hook V^*}(\cA)\succ0.
$$
\end{proposition}

\begin{proof}
The decomposition \eqref{eq:spin-decomposition} is orthogonal and invariant.
Positivity on $V_{4,2}$, $V_{2,4}$ and $V_{2,2}$ follows from
\eqref{eq:Qsplit} and Lemma \ref{lem:FT}. Positivity on $V_{2,0}$ and $V_{0,2}$
follows from \eqref{eq:two-form-positive}. Therefore the operator is positive
on every irreducible summand.
\end{proof}

\subsubsection{The 6-dimensional cubic estimate}\label{sec:six-dimensional}

In dimension $6$ the stable orthogonal Pieri calculation of
Section \ref{sec:tracefree-hook-weitzenbock} has low-rank coincidences.  It is simpler to retain
the full cubic space satisfying the normality constraint in the estimate of
Section \ref{sec:reduction}.

\begin{lemma}\label{lem:six-dimensional-cubic}
Let $n=6$, $k=3$, and let
$K\in\Sym^3_0V^*\otimes V$ satisfy
\eqref{eq:highest-normal-condition}. If
$\frac14\leq\sec_{\cA}\leq1$, then
\begin{equation}\label{eq:six-dimensional-gap}
  W_{\cA}(K)\geq\frac{27}{176}|K|^2.
\end{equation}
Consequently, a nonzero top cubic Fourier coefficient satisfying the
normality constraint and the twisted conformal Killing equation cannot occur on a closed $\delta$-pinched $6$-manifold with
$\delta>1/4$.
\end{lemma}

\begin{proof}
The derivation of \eqref{eq:Gnk} is valid for $k=3$; only the positivity
argument using $\Pi_{n,k}$ required $k\geq5$.  For $(n,k)=(6,3)$,
\eqref{eq:NPY-decomposition} gives
$0\leq u=N_{\parallel}/N\leq3/4$, and
$$
 G_{6,3}(u)=\frac{77}{16}-\frac{11}{4}u-\sqrt{21(1-u)}.
$$
Writing $z=\sqrt{1-u}$ gives the exact square
$$
 G_{6,3}(u)=\frac{27}{176}
 +\frac{11}{4}\left(z-\frac{2\sqrt{21}}{11}\right)^2.
$$
The minimiser $u=37/121$ lies in $[0,3/4]$, proving
\eqref{eq:six-dimensional-gap}. The final assertion follows from
\eqref{eq:tangent-twist-pestov-normalized}.
\end{proof}

\subsubsection{The stable cubic vanishing statement}

\begin{lemma}\label{lem:two-form-general}
Let $n$ be even and $0\leq\sec_{\cE}\leq1$.  On $\Lambda^2V^*$,
\begin{equation}\label{eq:two-form-general}
  \cQ_{\Lambda^2}(\cE)
  \succeq-\frac{2(n-2)}3\Id.
\end{equation}
Consequently, if $\delta>1/4$ and
$\cA=\delta I+(1-\delta)\cE$, then
\begin{equation}\label{eq:strict-two-form-general}
 \cQ_{\Lambda^2}(\cA)
 \succeq\frac{2(n-2)}3(4\delta-1)\Id\succ0.
\end{equation}
\end{lemma}

\begin{proof}
Write
$\omega=\sum_{a=1}^{n/2}\lambda_a e_{2a-1}\wedge e_{2a}$.
For $a<b$, put
$$
 \Theta_{ab}:=\sum_{u\in\{2a-1,2a\}}
              \sum_{v\in\{2b-1,2b\}}\sec_{\cE}(e_u,e_v),
 \qquad
 r_{ab}:=\cE_{2a-1,2a,2b-1,2b}.
$$
The two-form curvature formula is
\begin{equation*}
 \inner{\cQ_{\Lambda^2}(\cE)\omega}{\omega}
 =\sum_{a<b}\left\{
 \Theta_{ab}(\lambda_a^2+\lambda_b^2)
 -4r_{ab}\lambda_a\lambda_b\right\}.
\end{equation*}
Here $\Theta_{ab}\geq0$, while Lemma~\ref{lem:BBK} gives
$|r_{ab}|\leq2/3$.  Hence each pair is bounded below by
$$
 -\frac83|\lambda_a\lambda_b|
 \geq-\frac43(\lambda_a^2+\lambda_b^2).
$$
After summing, every $\lambda_a^2$ occurs $n/2-1$ times, proving
\eqref{eq:two-form-general}.  Since
$\cQ_{\Lambda^2}(I)=2(n-2)\Id$, linearity gives
\eqref{eq:strict-two-form-general}.
\end{proof}

\begin{proposition}\label{prop:stable-cubic}
Let $n\geq10$ be even, let $(M^n,g)$ be closed and $\delta$-pinched for some
$\delta>1/4$, and let
$T\in C^\infty(M,\cC)$ satisfy $\cD T=0$. Then $T=0$.
\end{proposition}

\begin{proof}
Let $\kappa_0$ be a pinching scale and let $\cA=\kappa_0^{-1}\Rop^g$ be the
normalized curvature operator. Write $T=H+S_0+\omega$ according to
\eqref{eq:orthogonal-hook-decomposition}. Since
$$
  \cA=\frac14I+\frac34\cE,
  \qquad 0\leq\sec_{\cE}\leq1,
$$
\eqref{eq:partial-integral-sign} and
\eqref{eq:tracefree-hook-twist-positive} give
$$
  0\geq\int_Mw_{\Rop^g}(H,H)\,d\vol_g
  =\kappa_0\int_Mw_{\cA}(H,H)\,d\vol_g
  \geq\kappa_0c_n\int_M|H|^2\,d\vol_g.
$$
Hence $H=0$.

Return to the ordinary formula of Appendix \ref{sec:weitzenbock}. Since $\cD T=0$,
Corollary \ref{cor:integral-sign} gives
\begin{equation}\label{eq:trace-block-integral}
  0\leq-\kappa_0\int_M
  \inner{\cQ_{\Hook}(\cA)T}{T}\,d\vol_g.
\end{equation}
The two remaining summands in
\eqref{eq:orthogonal-hook-decomposition} are invariant under every
infinitesimal orthogonal action, so the curvature term has no cross term
between them. For $S_0\in\Sym^2_0V^*$, diagonalized with eigenvalues
$\lambda_i$,
\begin{equation*}
  \inner{\cQ_{\Sym^2}(\cA)S_0}{S_0}
  =2\sum_{i<j}\sec_{\cA}(e_i,e_j)(\lambda_i-\lambda_j)^2>0
\end{equation*}
unless $S_0=0$.

Choose $\delta_1$ with $1/4<\delta_1<\delta$. Then
$$
  \cA=\delta_1I+(1-\delta_1)\cE_0,
  \qquad 0\leq\sec_{\cE_0}\leq1.
$$
By Lemma \ref{lem:two-form-general}, the curvature term on $\Lambda^2V^*$ is
strictly positive. Therefore the right-hand side of
\eqref{eq:trace-block-integral} is non-positive and can vanish only when
$S_0=\omega=0$. Thus $T=0$.
\end{proof}

\section{The invariant projector obstruction}\label{sec:symmetric-projector}

We now treat the obstruction furnished by an invariant orthogonal projector.
Throughout this Section, integration over a Euclidean unit sphere is taken
with respect to normalized spherical measure.
\subsection{The symmetric-square Pestov curvature form}

Let $V$ be an $n$-dimensional Euclidean space and let
$$
  S:S(V)\longrightarrow\Sym^2V
$$
be a smooth map. We regard $S(v)$ as a self-adjoint endomorphism of $V$.
For a fixed orthonormal basis $(e_a)_{a=0}^{n-1}$, held fixed during
spherical differentiation, put
$$
  S_{ab}(v):=\inner{S(v)e_b}{e_a},
  \qquad
  z_{ab}:=\nabla^S S_{ab},
$$
and
$$
  \beta_{ab}:=v\wedge z_{ab},
  \qquad
  \alpha_{ab}:=S(v)e_b\wedge e_a.
$$
For an algebraic curvature operator $\cB$ on $V$, define
\begin{equation}\label{eq:sym2-fibrewise-pestov-form}
  \mathscr W^{(2)}_{\cB}(S)
  :=\int_{S(V)}\sum_{a,b=0}^{n-1}
  \left\{
    \cB(\beta_{ab},\beta_{ab})
    -2\cB(\alpha_{ab},\beta_{ab})
  \right\}\,d\sigma(v).
\end{equation}
The sum is independent of the fixed orthonormal basis. When $S$ is the
restriction of a homogeneous harmonic polynomial of degree $k$, we use the same
notation for this degree-$k$ coefficient. The definition is made for arbitrary
smooth $S$; in particular, for
$P\in C^\infty(\SM,\Sym^2\cN)$, the expression
$\mathscr W^{(2)}_{\cB}(P)$ is obtained fibrewise from the full map
$v\mapsto P(x,v)$, without first separating its Fourier modes.

\begin{lemma}
\label{lem:sym2-pestov-sign}
Let $S$ be the highest Fourier coefficient of a flow-invariant section of
$\Sym^2\cN$. Then
\begin{equation}\label{eq:sym2-localized-sign}
  \int_M\mathscr W^{(2)}_{\Rop^g}(S)\,d\vol_g\leq0.
\end{equation}
If $P\in C^\infty(\SM,\Sym^2\cN)$ itself satisfies $XP=0$, then
\begin{equation}\label{eq:sym2-full-sign}
  \int_M\mathscr W^{(2)}_{\Rop^g}(P)\,d\vol_g\leq0.
\end{equation}
The same statements hold with $\Rop^g$ replaced by the normalized operator
$\cA=\kappa_0^{-1}\Rop^g$.
\end{lemma}

\begin{proof}
For the induced connection on $\Sym^2TM$, the curvature acts by
$$
  R^{\Sym^2}(u,w)S=[R^g(u,w),S].
$$
If $D_iS:=\nabla^V_{e_i}S$, then, using that $S$ and $D_iS$ are
self-adjoint and $R^g(v,e_i)$ is skew-adjoint,
\begin{align*}
 \sum_i\inner{[R^g(v,e_i),S]}{D_iS}
 &=2\sum_{i,a,b}
   \inner{R^g(v,e_i)S e_b}{e_a}(D_iS)_{ab} \\
 &=2\sum_{a,b}\Rop^g(\alpha_{ab},\beta_{ab}).
\end{align*}
The Jacobi term in the twisted Pestov identity is
$-\sum_{a,b}\Rop^g(\beta_{ab},\beta_{ab})$.  Thus
\eqref{eq:sym2-fibrewise-pestov-form} is exactly the fibrewise integrand of the symmetric-square Pestov curvature form
\eqref{eq:sym2-pestov-form-normalized}.  The localized sign
\eqref{eq:sym2-localized-sign} follows from \eqref{eq:top-pestov-sign}, and the
full sign \eqref{eq:sym2-full-sign} follows from \eqref{eq:full-pestov-sign}.
Multiplication of the curvature operator by the positive scalar
$\kappa_0^{-1}$ does not change either sign.
\end{proof}

For a highest even Fourier coefficient coming from an orthogonal projector,
\cite[Lemma~4.2]{CLMS24} gives the normal conditions
\begin{equation}\label{eq:sym2-normal-conditions}
  \tr S=0,
  \qquad
  S(v)v\ \text{has degree }k-1,
  \qquad
  \inner{S(v)v}{v}\ \text{has degree }k-2.
\end{equation}
The first identity also follows directly from the fact that the trace of the
projector is constant.

Put
\begin{equation*}
  N:=\int_{S(V)}|S|^2,
  \qquad
  N_1:=\int_{S(V)}|S(v)v|^2,
  \qquad
  \lambda_k:=k(n+k-2).
\end{equation*}

\begin{lemma}[General estimate for the symmetric-square Pestov form]
\label{lem:sym2-general-estimate}
Assume that $S$ satisfies the first two conditions in
\eqref{eq:sym2-normal-conditions}. Then
\begin{equation}\label{eq:sym2-WI}
  \mathscr W^{(2)}_I(S)
  =(\lambda_k-2)N-2(n+2k-4)N_1.
\end{equation}
If $0\leq\sec_{\cE}\leq1$, then
\begin{equation}\label{eq:sym2-residual}
 \mathscr W^{(2)}_{\cE}(S)
 \geq
 -N-(n-2)N_1
 -\frac43\sqrt{(n-2)(N-N_1)\lambda_kN}.
\end{equation}
\end{lemma}

\begin{proof}
For $I$, the first term in \eqref{eq:sym2-fibrewise-pestov-form} is
$\lambda_kN$.  The constant-curvature computation for the commutator action
on $\Sym^2V$ is
$$
  2N+2(n+2k-4)N_1;
$$
this is exactly the symmetric-square identity in
\cite[Lemma~4.6]{CLMS24}. Subtraction gives
\eqref{eq:sym2-WI}.

We prove \eqref{eq:sym2-residual} directly. Fix $v\in S(V)$, put
$e_0=v$, and write
$$
  S(v)=
  \begin{pmatrix}
    a&b^\top\\ b&C
  \end{pmatrix}
$$
relative to $V=\bbR v\oplus v^\perp$.  If
$S(v)e_j=s_{0j}v+y_j$ with $y_j\perp v$, then the radial part of
$\alpha_{0j}$ is $-v\wedge y_j$, whereas the radial part of
$\alpha_{ij}$ is $s_{0j}v\wedge e_i$ for $i\geq1$.  The Jacobi form
$$
  \mathsf J_v^{\cE}(u,w):=\cE(v\wedge u,v\wedge w)
$$
satisfies $0\preceq\mathsf J_v^{\cE}\preceq\Id$. Hence
$$
  \mathsf J_v^{\cE}(z,z)-2\mathsf J_v^{\cE}(q,z)
  =\mathsf J_v^{\cE}(z-q,z-q)-\mathsf J_v^{\cE}(q,q)
  \geq-|q|^2.
$$
After summing all radial terms, the total loss is
$$
  \sum_j|y_j|^2+(n-1)\sum_j s_{0j}^2
  =|S|^2+(n-2)|S(v)v|^2.
$$

The remaining terms are
$$
 -2\sum_{i=1}^{n-1}\sum_{j=0}^{n-1}
 \cE(y_j\wedge e_i,v\wedge z_{ij}).
$$
The two displayed simple forms are orthogonal.  Lemma~\ref{lem:BBK} gives
$$
  |\cE(y_j\wedge e_i,v\wedge z_{ij})|
  \leq\frac23|y_j\wedge e_i|\,|z_{ij}|.
$$
Moreover,
$$
  \sum_{i,j}|y_j\wedge e_i|^2
  =(n-2)\sum_j|y_j|^2
  =(n-2)(|S|^2-|S(v)v|^2).
$$
Cauchy--Schwarz in $(v,i,j)$ and
$$
  \int_{S(V)}\sum_{a,b}|z_{ab}|^2
  =\int_{S(V)}|\nabla^SS|^2
  =\lambda_kN
$$
now prove \eqref{eq:sym2-residual}.
\end{proof}

\subsection{Fourier degrees at least six}

\begin{proposition}
\label{prop:sym2-high-degree}
Let $n\geq12$ and $k\geq6$ be even. If
$$
  \frac5{13}\leq\sec_{\cA}\leq1,
$$
then every nonzero $S$ satisfying \eqref{eq:sym2-normal-conditions} obeys
$$
  \mathscr W^{(2)}_{\cA}(S)>0.
$$
\end{proposition}

\begin{proof}
Write
$$
  \cA=\frac5{13}I+\frac8{13}\cE,
  \qquad 0\leq\sec_{\cE}\leq1.
$$
Let $u=N_1/N$, $z=\sqrt{1-u}$, and put
$$
  Q:=(n-2)\lambda_k,
  \qquad
  L:=10k+9n-28.
$$
By Lemma \ref{lem:sym2-general-estimate},
\begin{align}
 \frac{13}{8N}\mathscr W^{(2)}_{\cA}(S)
 &\geq
 \frac58\{\lambda_k-2-2(n+2k-4)u\}
 -1-(n-2)u-\frac43\sqrt{Q(1-u)}\notag\\
 &=\frac L4z^2-\frac43\sqrt Q\,z
   +\frac{5k^2+5kn-30k-18n+38}{8}\notag\\
 &=\frac L4\left(z-\frac{8\sqrt Q}{3L}\right)^2
   +\frac{\Psi_{n,k}}{72L},
 \label{eq:sym2-high-square}
\end{align}
where
\begin{align}
 \Psi_{n,k}={}&450k^3+727k^2n-3704k^2
 +277kn^2-4798kn\notag\\
 &+10468k-1458n^2+7614n-9576.
 \label{eq:sym2-Psi}
\end{align}
Writing $n=12+a$ and $k=6+b$ gives
\begin{align*}
 \Psi_{12+a,6+b}={}&277a^2b+204a^2+727ab^2+10574ab+9894a\\
 &+450b^3+13120b^2+101620b+106440.
\end{align*}
Every coefficient is positive. Thus \eqref{eq:sym2-high-square} is strictly
positive whenever $S\neq0$.
\end{proof}

\subsection{The quartic mode}

For an admissible quartic coefficient, fix $v\in S(V)$, put $e_0=v$, and write
$$
  S(v)=\begin{pmatrix}a&b^\top\\b&C\end{pmatrix}.
$$
With $z_{ab}=\nabla^SS_{ab}$, define the following invariant quadratic
forms:
\begin{align}
 A_0&:=\int\bigl((n-1)a^2+n|b|^2+|C|^2\bigr),
 &B&:=\int|b|^2,\notag\\
 A_C&:=\int\bigl((n-1)|C|^2-a^2\bigr),
 &Z_b&:=\int\sum_{\alpha=1}^{n-1}|z_{\alpha0}|^2,\notag\\
 &&Z_C&:=\frac12\int\sum_{\alpha,\beta=1}^{n-1}
 |z_{\alpha\beta}|^2.\notag
\end{align}
Here and below Greek indices range from $1$ to $n-1$.

\begin{lemma}
\label{lem:quartic-residual}
If $0\leq\sec_{\cE}\leq1$, then
\begin{equation}\label{eq:quartic-residual}
 \mathscr W^{(2)}_{\cE}(S)
 \geq-A_0-\frac43
 \left(\sqrt{A_CZ_C}+\sqrt{(n-2)BZ_b}\right).
\end{equation}
\end{lemma}

\begin{proof}
The completion of squares in the proof of
Lemma \ref{lem:sym2-general-estimate} gives the radial loss $-A_0$.  For the
remaining terms, diagonalize $C$ at the chosen point, say
$Ce_\alpha=\lambda_\alpha e_\alpha$. Since $z_{\alpha\beta}=z_{\beta\alpha}$,
\begin{align*}
 &\sum_{\alpha,\beta=1}^{n-1}
 \cE(Ce_\beta\wedge e_\alpha,v\wedge z_{\alpha\beta})=
 \sum_{\alpha<\beta}(\lambda_\alpha-\lambda_\beta)
 \cE(e_\alpha\wedge e_\beta,v\wedge z_{\alpha\beta}).
\end{align*}
The admissible coefficient is trace-free, so $\tr C=-a$, and therefore
$$
 \sum_{\alpha<\beta}(\lambda_\alpha-\lambda_\beta)^2
 =(n-1)|C|^2-a^2.
$$
The sum of the squares of the corresponding off-diagonal
$z_{\alpha\beta}$ is at most the integrand defining $Z_C$.  Lemma~\ref{lem:BBK}
and Cauchy--Schwarz thus bound this part of the mixed
term by $\frac23\sqrt{A_CZ_C}$.

The terms involving $b$ are
$$
 \sum_\alpha\cE(b\wedge e_\alpha,v\wedge z_{\alpha0}).
$$
Since $\sum_\alpha|b\wedge e_\alpha|^2=(n-2)|b|^2$, their absolute value is
at most $\frac23\sqrt{(n-2)BZ_b}$.  The coefficient $2$ in
\eqref{eq:sym2-fibrewise-pestov-form} produces the factor $4/3$ in
\eqref{eq:quartic-residual}.
\end{proof}

\begin{proposition}
\label{prop:sym2-quartic}
Let $n\geq12$. If
$\frac5{13}\leq\sec_{\cA}\leq1$, then
$
  \mathscr W^{(2)}_{\cA}(S)>0
$
for every nonzero $S\in\mathscr N^{(2)}_4$.
\end{proposition}

\begin{proof}
Write $\cA=\frac5{13}I+\frac8{13}\cE$.  From
Lemma \ref{lem:quartic-residual} and the two Young inequalities
\begin{align}
 \frac43\sqrt{A_CZ_C}
 &\leq2A_C+\frac29Z_C,\label{eq:quartic-young-one}\\
 \frac43\sqrt{(n-2)BZ_b}
 &\leq\frac{16}{9}(n-2)B+\frac14Z_b,
 \notag
\end{align}
we obtain
\begin{align}
 \frac{13}{8}\mathscr W^{(2)}_{\cA}(S)
 \geq{}&\frac58\mathscr W_I^{(2)}-A_0-2A_C-\frac29Z_C\notag\\
 &-\frac{16}{9}(n-2)B-\frac14Z_b.
 \label{eq:quartic-master-gap}
\end{align}
All terms on the right are $O(n)$-invariant quadratic forms.  Since
\eqref{eq:quartic-decomposition} is multiplicity-free, there are no cross
terms between its three summands. Substitution of
Table \eqref{eq:quartic-table} gives the respective gaps
\begin{equation}\label{eq:quartic-gaps}
 \frac{n+69}{18},
 \qquad
 \frac{43n+220}{9(n+2)},
 \qquad
 \frac{2(8n^2+67n+110)}{9n(n+1)}.
\end{equation}
Every number in \eqref{eq:quartic-gaps} is strictly positive. Therefore the
right-hand side of \eqref{eq:quartic-master-gap} is positive unless all
three components of $S$ vanish.
\end{proof}

\subsection{The quadratic mode}

We finish the symmetric-projector branch by exploiting the additional
algebraic symmetry of a degree-two projector.

\begin{lemma}
\label{lem:quadratic-young}
Let $Q:V\to\Sym^2V$ be a homogeneous quadratic polynomial satisfying
$Q(x)x=0$ for every $x$. Write $Q(x)=\mathsf B(x,x)$, where
$\mathsf B$ is symmetric bilinear with values in $\Sym^2V$, and set
$$
 \mathsf A(u,w,y,z):=\inner{\mathsf B(u,w)y}{z}.
$$
Then
\begin{equation}\label{eq:quadratic-young-relations}
 \mathsf A(u,w,y,z)=\mathsf A(w,u,y,z)
 =\mathsf A(u,w,z,y),
 \qquad
 \mathsf A_{(uwy)z}=0,
\end{equation}
and, in addition,
\begin{equation}\label{eq:quadratic-pair-symmetry}
 \mathsf A(u,w,y,z)=\mathsf A(y,z,u,w).
\end{equation}
\end{lemma}

\begin{proof}
The first two symmetries are immediate. Polarizing
$\mathsf A(x,x,x,z)=0$ gives the Young relation in
\eqref{eq:quadratic-young-relations}.  Write
$$
 Y(u,w,y;z):=\mathsf A(u,w,y,z)+\mathsf A(w,y,u,z)
              +\mathsf A(y,u,w,z).
$$
Thus $Y=0$.  Adding $Y(u,w,y;z)=0$ and $Y(u,w,z;y)=0$ gives
$2\mathsf A(u,w,y,z)$ plus four further terms.  Adding
$Y(y,z,u;w)=0$ and $Y(y,z,w;u)=0$ gives
$2\mathsf A(y,z,u,w)$ plus the same four terms, in a different order, by
the within-pair symmetries.  Subtracting the two identities proves
\eqref{eq:quadratic-pair-symmetry}.
\end{proof}

\begin{proposition}[Elimination of degree-two projectors]
\label{prop:quadratic-projector}
Let $n$ be even, let $(M^n,g)$ be closed, and suppose that
$\frac5{13}\leq\sec_{\cA}\leq1$. There is no flow-invariant orthogonal
projector
$$
 P\in C^\infty(\SM,\Sym^2\cN)
$$
of constant rank
$$
  1\leq r\leq n-2
$$
and vertical Fourier degree at most two.
\end{proposition}

\begin{proof}
Suppose that such a projector $P$ exists, and denote its rank by $r$.
Fix $x\in M$ and carry out the following algebra on $V=T_xM$.  By
\cite[Lemma~4.2]{CLMS24},
\begin{equation*}
  P=P_0+P_2,
  \qquad
  P_0=\frac rn\Id_V.
\end{equation*}
For completeness, the last identity follows by representing $P$ by a
homogeneous quadratic tensor and using the polarized normal relation
$K(v,v,w,w)=K(w,w,v,v)$; tracing in $w$ gives
$r=n\inner{P_0v}{v}$.

Fix $v\in S(V)$ and decompose
$$
 V=\bbR v\oplus E_v\oplus F_v,
 \qquad
 E_v=\operatorname{Im}P(v),
 \qquad
 \dim E_v=r,
 \qquad
 \dim F_v=s:=n-1-r.
$$
The rank assumption gives $r\geq1$ and $s\geq1$. Choose an adapted
orthonormal basis
$e_0=v$, $(e_\alpha)$ of $E_v$, and $(e_p)$ of $F_v$.  Differentiating
$P^2=P$ shows that $\nabla^V_wP$ has only off-diagonal blocks between
$E_v$ and $\ker P(v)=\bbR v\oplus F_v$. Differentiating $P(v)v=0$ gives
$$
  z_{\alpha0}=-e_\alpha,
  \qquad
  z_{p0}=0.
$$
The only free components are therefore $z_{\alpha p}=z_{p\alpha}$.  In
\eqref{eq:sym2-fibrewise-pestov-form}, the two quadratic terms associated with
$z_{\alpha0}$ and $z_{0\alpha}$ are exactly cancelled by the twist term
with indices $(0,\alpha)$.  The two copies of each free quadratic term
remain, while only the indices $(p,\alpha)$ contribute a free twist term.
Consequently,
\begin{equation}\label{eq:projector-exact-W}
 \mathscr W^{(2)}_{\cB}(P)
 =2\int_{S(V)}\sum_{\alpha,p}
 \left\{
  \cB(v\wedge z_{\alpha p},v\wedge z_{\alpha p})
  -\cB(e_\alpha\wedge e_p,v\wedge z_{\alpha p})
 \right\}.
\end{equation}

The degree-two energy is exact.  Since $\lambda_2=2n$,
\begin{align*}
 \int|\nabla^VP|^2
 &=2n\int|P_2|^2
 =2n\left(r-\frac{r^2}{n}\right)
 =2r(n-r).
\end{align*}
On the other hand, the forced components contribute $2r$ and the free
components occur twice. Hence, with
\begin{equation*}
 Z:=\int\sum_{\alpha,p}|z_{\alpha p}|^2,
\end{equation*}
we have
\begin{equation*}
  Z=r(n-1-r)=rs.
\end{equation*}

It remains to improve the estimate of the mixed term.  Define the quadratic
homogeneous extension
$$
  Q(x):=|x|^2P\left(\frac{x}{|x|}\right).
$$
Lemma \ref{lem:quadratic-young} applies to $Q$.  Set
$$
 T_{\alpha\beta p}:=\inner{z_{\alpha p}}{e_\beta},
 \qquad
 U_{\alpha pq}:=\inner{z_{\alpha p}}{e_q}.
$$
The Young relation, pair symmetry, and the off-diagonal property of
$\nabla^VP$ imply
\begin{equation*}
  T_{\alpha\beta p}=-T_{\beta\alpha p},
  \qquad
  U_{\alpha pq}=-U_{\alpha qp}.
\end{equation*}
Indeed, for example, the polarized identity applied to
$(v,e_\beta,e_p;e_\alpha)$ has first two nonzero terms
$\frac12T_{\alpha\beta p}$ and
$\frac12T_{\beta\alpha p}$, while its third term is an $E_v$--$E_v$ block
of a derivative of $P$ and vanishes.  The proof for $U$ is identical, using
the vanishing $F_v$--$F_v$ block.

Put
$$
 \tau_p:=\sum_{\alpha<\beta}T_{\alpha\beta p}
 e_\alpha\wedge e_\beta,
 \qquad
 \upsilon_\alpha:=\sum_{p<q}U_{\alpha pq}e_p\wedge e_q.
$$
The first Bianchi identity rewrites the mixed term as
\begin{equation*}
 M_{\cE}:=\sum_{\alpha,p}
 \cE(e_\alpha\wedge e_p,v\wedge z_{\alpha p})
 =\sum_\alpha\cE(e_\alpha\wedge v,\upsilon_\alpha)
  +\sum_p\cE(\tau_p,v\wedge e_p).
\end{equation*}
Putting each two-form in canonical form and applying Lemma~\ref{lem:BBK}
to every simple summand yields
\begin{equation*}
 |M_{\cE}|
 \leq\frac23\left(
 \sqrt{\left\lfloor\frac s2\right\rfloor}
 \sum_\alpha|\upsilon_\alpha|
 +\sqrt{\left\lfloor\frac r2\right\rfloor}
 \sum_p|\tau_p|
 \right).
\end{equation*}
Moreover,
\begin{equation}\label{eq:projector-TU-energy}
 \sum_{\alpha,p}|z_{\alpha p}|^2
 =2\sum_p|\tau_p|^2+2\sum_\alpha|\upsilon_\alpha|^2.
\end{equation}
Two applications of Cauchy--Schwarz, first in the indices and then between
the two summands in \eqref{eq:projector-TU-energy}, give
\begin{equation*}
 \int|M_{\cE}|
 \leq C_{r,s}Z,
 \qquad
 C_{r,s}:=\frac{2}{3\sqrt2}
 \sqrt{
  \frac{\lfloor r/2\rfloor}{r}
  +\frac{\lfloor s/2\rfloor}{s}
 }.
\end{equation*}
Since $r+s=n-1$ is odd, $r$ and $s$ cannot both be even. Therefore
\begin{equation*}
  C_{r,s}<\frac{\sqrt2}{3}<\frac58.
\end{equation*}

Finally write $\cA=\frac5{13}I+\frac8{13}\cE$.  In
\eqref{eq:projector-exact-W}, the $I$-mixed term vanishes because a
tangential two-form is orthogonal to a radial two-form, while the first
$\cE$-term is nonnegative.  Consequently,
\begin{align*}
 \mathscr W^{(2)}_{\cA}(P)
 &\geq\frac{16}{13}
 \left(\frac58Z-\int|M_{\cE}|\right)\\
 &\geq\frac{16}{13}\left(\frac58-C_{r,s}\right)Z>0.
\end{align*}
The estimate is fibrewise in $x$ and is strict because $Z=rs>0$.
Integrating over $M$ contradicts the full Pestov sign
\eqref{eq:sym2-full-sign}.
\end{proof}

\begin{proposition}
\label{prop:symmetric-projector-vanishing}
Let $n\geq12$ be divisible by $4$ and let $(M^n,g)$ be closed and
$5/13$-pinched. Then there is no nontrivial (that is, neither $0$ nor
$\Id_{\cN}$) even flow-invariant orthogonal projector in
$C^\infty(\SM,\Sym^2\cN)$.
\end{proposition}

\begin{proof}
Let $P$ be such a projector. It has finite even Fourier degree. If its
highest degree is $k\geq6$, Proposition \ref{prop:sym2-high-degree} and the
localized Pestov sign give a contradiction. If the highest degree is $4$,
Proposition \ref{prop:sym2-quartic} gives the same contradiction. Hence
$\deg P\leq2$. Degree zero is impossible because an endomorphism independent
of $v$ and annihilating every $v$ must vanish. Thus $P$ has degree two. Since
$P$ is neither $0$ nor $\Id_{\cN}$, its rank satisfies $1\leq r\leq n-2$,
and Proposition \ref{prop:quadratic-projector} gives the final contradiction.
\end{proof}

\section{Exceptional form-valued obstructions}\label{sec:exceptional-dimensions}

The form-valued obstructions in dimensions $8$ and $134$ are most
efficiently treated without separating their highest Fourier coefficients.
We first establish a Weitzenb\"ock sign for the spherical evaluation inner
product and the resulting orbit-type curvature identity.  The $7$-dimensional obstruction is of a
different type and is treated at the end of the section by a short cubic
refinement.

\subsection{The spherical evaluation inner product and the orbit-type curvature identity}
\label{sec:orbit-type-framework}

Let $V$ be an $n$-dimensional Euclidean vector space, and for integers
$k\geq1$ and $1\leq p\leq n-1$ put
\begin{equation}\label{eq:normal-form-module}
 E_{k,p}(V):=
 \left\{
 T\in\Sym^kV^*\otimes\Lambda^pV^*:
 \iota_vT(v,\ldots,v)=0\ \text{for every }v\in V
 \right\}.
\end{equation}
This is an $\mathrm O(V)$-module.  On it we use the spherical evaluation inner
product
\begin{equation}\label{eq:evaluation-inner-product}
 \inner{S}{T}_{\mathrm{ev}}
 :=\int_{S(V)}
 \inner{S(v,\ldots,v)}{T(v,\ldots,v)}_{\Lambda^pV^*}
 \,d\sigma(v),
\end{equation}
where $d\sigma$ is normalized spherical measure.  This inner product is
positive definite and $\mathrm O(V)$-invariant.  We use the induced product
on $V^*\otimes E_{k,p}(V)$.

Let $\rho$ be the infinitesimal orthogonal action on $E_{k,p}(V)$ and define
\begin{equation}\label{eq:exceptional-B}
 B_{k,p}(a\otimes T)
 :=\sum_i e^i\otimes\rho(e_i\wedge a)T.
\end{equation}
The generalized-Killing symbol is
\begin{equation*}
 \sigma_{k,p}(U)(v^{k+1})=U_v(v^k),
 \qquad U\in V^*\otimes E_{k,p}(V).
\end{equation*}

\begin{lemma}
\label{lem:exceptional-symbol-sign}
If $U\in V^*\otimes E_{k,p}(V)$ satisfies $\sigma_{k,p}(U)=0$, then
\begin{equation}\label{eq:exceptional-symbol-sign}
 \inner{B_{k,p}U}{U}_{\mathrm{ev}}\leq0.
\end{equation}
More precisely, for a unit vector $v$, put $W=v^\perp$, choose an
orthonormal basis $e_0=v,e_1,\ldots,e_{n-1}$, write
$U=\sum_i e^i\otimes U_i$, and set
$$
 u_a:=U_a(v^k)\in\Lambda^pW^*,
 \qquad
 \eta_v:=\sum_{a=1}^{n-1}e^a\otimes u_a
 \in W^*\otimes\Lambda^pW^*.
$$
Then the integrand of the left-hand side of
\eqref{eq:exceptional-symbol-sign} is
\begin{equation*}
 -|\varepsilon\eta_v|^2-|\iota\eta_v|^2,
 \qquad
 \varepsilon\eta_v:=\sum_a e^a\wedge u_a,
 \quad
 \iota\eta_v:=\sum_a\iota_{e_a}u_a.
\end{equation*}
\end{lemma}

\begin{proof}
The normal relation in \eqref{eq:normal-form-module} gives
$U_i(v^k)\in\Lambda^pW^*$ for every $i$.  At $v=e_0$, the symbol equation
first gives
\begin{equation*}
 u_0:=U_0(v^k)=0.
\end{equation*}
Differentiate the polynomial identity $U_y(y^k)=0$ at $y=v$ in the
direction $e_a\in W$.  This gives
\begin{equation*}
 u_a+kU_0(e_a,v^{k-1})=0.
\end{equation*}
With the convention \eqref{eq:wedge-skew-convention}, one has
$(e_a\wedge e_0)v=-e_a$.  Hence
\begin{equation*}
 \bigl(\rho(e_a\wedge e_0)U_0\bigr)(v^k)
 =kU_0(e_a,v^{k-1})=-u_a.
\end{equation*}
For $a,b\geq1$, the skew endomorphism $e_a\wedge e_b$ fixes $v$, so its
action on $U_b(v^k)$ is only the ordinary action on
$\Lambda^pW^*$.  It follows from \eqref{eq:exceptional-B} that
\begin{align}
 \sum_i\inner{(B_{k,p}U)_i(v^k)}{u_i}
 &=-\sum_a|u_a|^2
   +\sum_{a,b}
    \inner{\rho_{\Lambda^pW}(e_a\wedge e_b)u_b}{u_a}.
 \label{eq:orbit-pointwise-B}
\end{align}
Let $B_W^{(p)}$ denote the conformal-weight operator on
$W^*\otimes\Lambda^pW^*$.  Direct expansion of exterior multiplication and
contraction gives the standard identity
\begin{equation*}
 \Id-B_W^{(p)}=\varepsilon^*\varepsilon+\iota^*\iota.
\end{equation*}
Thus the right-hand side of \eqref{eq:orbit-pointwise-B} is exactly
$-|\varepsilon\eta_v|^2-|\iota\eta_v|^2$.  Integrating over $S(V)$ proves
the assertion.
\end{proof}

Let $\mathscr E_{k,p}\to M$ be the associated bundle whose fibre at $x$
is $E_{k,p}(T_xM)$.  For a smooth section $T$ define
\begin{equation}\label{eq:exceptional-Killing-operator}
 (\cD_{k,p}T)(v^{k+1}):=(\nabla_vT)(v^k).
\end{equation}
The expression on the right is homogeneous of degree $k+1$ in $v$, so
\eqref{eq:exceptional-Killing-operator} determines the full symmetrized
covariant derivative.

\begin{corollary}
\label{cor:exceptional-integral-sign}
Let $(M^n,g)$ be closed and let
$T\in C^\infty(M,\mathscr E_{k,p})$ satisfy $\cD_{k,p}T=0$.  If $g$ is
$\delta$-pinched at scale $\kappa_0$ and
$\cA=\kappa_0^{-1}\Rop^g$, then
\begin{equation*}
 \int_M
 \inner{\cQ_{\mathscr E_{k,p}}(\cA)T}{T}_{\mathrm{ev}}
 \,d\vol_g\leq0.
\end{equation*}
\end{corollary}

\begin{proof}
The inner product \eqref{eq:evaluation-inner-product} is
$\mathrm O(n)$-invariant and hence parallel on the associated bundle.  In
particular every infinitesimal orthogonal action is skew-adjoint, and
$B_{k,p}$ is self-adjoint.  The universal identity, with the type-correct
contraction introduced in Appendix~\ref{sec:weitzenbock}, is
$$
 \widetilde B_{k,p}(\nabla^2T)=q_{\mathscr E_{k,p}}(R^g)T.
$$
Integration by parts and Lemma~\ref{lem:exceptional-symbol-sign} give
$$
 \int_M\inner{q_{\mathscr E_{k,p}}(R^g)T}{T}_{\mathrm{ev}}
 =-\int_M\inner{B_{k,p}\nabla T}{\nabla T}_{\mathrm{ev}}\geq0.
$$
Now use $q_{\mathscr E_{k,p}}(R^g)=-\kappa_0\cQ_{\mathscr E_{k,p}}(\cA)$.
\end{proof}

\begin{lemma}[Homogeneous tensor associated with a finite Fourier expansion]
\label{lem:finite-fourier-homogenization}
Let
$$
 \varphi\in C^\infty(\SM,\Lambda^p\cN),
 \qquad X\varphi=0,
$$
have finite Fourier degree and pure parity, and let $k$ be its highest
Fourier degree.  If $\widehat\varphi_j$ denotes the homogeneous harmonic
extension of its degree-$j$ coefficient, then
\begin{equation}\label{eq:finite-fourier-homogenization}
 F_x(y):=
 \sum_{\substack{j\leq k\\j\equiv k\,\mathrm{mod}\,2}}
 |y|^{k-j}\widehat\varphi_{j,x}(y)
\end{equation}
is a homogeneous polynomial of degree $k$.  Polarizing $F$ defines
\begin{equation*}
 T\in C^\infty
 \bigl(M,\Sym^kT^*M\otimes\Lambda^pT^*M\bigr),
 \qquad T_x(v^k)=F_x(v),
\end{equation*}
with
$$
 T\in C^\infty(M,\mathscr E_{k,p}),
 \qquad \cD_{k,p}T=0.
$$
\end{lemma}

\begin{proof}
The powers $k-j$ are even, so \eqref{eq:finite-fourier-homogenization} is
polynomial and homogeneous of degree $k$.  On $S_xM$ one has
$F_x(v)=\varphi(x,v)$.  The normality relation
$\iota_v\varphi(x,v)=0$ therefore gives
$T\in\mathscr E_{k,p}$.  Along a geodesic, $|v|$ is constant and $v$ is
parallel, so $X\varphi=0$ is equivalent to
$(\nabla_vT)(v^k)=0$.  This is precisely $\cD_{k,p}T=0$.
\end{proof}

We now isolate the common curvature calculation.  Assume in addition that
the values of $\varphi$ have one fixed $\SO(n-1)$-orbit type.  Let $T$ be supplied by
Lemma~\ref{lem:finite-fourier-homogenization}.  Fix $(x,v)\in\SM$, put
$W=v^\perp$, and abbreviate
$$
 \phi:=T_x(v^k)=\varphi(x,v)\in\Lambda^pW^*.
$$
Let $H\subset\SO(W)$ be the stabilizer of $\phi$, let
$\mathfrak h\subset\Lambda^2W$ be its Lie algebra, and put
\begin{equation*}
 \Lambda^2W=\mathfrak h\oplus\mathfrak m,
 \qquad \mathfrak m:=\mathfrak h^\perp.
\end{equation*}
The orbit differential is
\begin{equation}\label{eq:orbit-differential}
 \Phi_\phi:\mathfrak m\longrightarrow\Lambda^pW^*,
 \qquad \Phi_\phi(Z)=\rho(Z)\phi.
\end{equation}
In the two applications below $\mathfrak m$ is irreducible, and hence
Schur's lemma gives a constant $\gamma>0$ such that
\begin{equation}\label{eq:orbit-homothety}
 \inner{\Phi_\phi(Z)}{\Phi_\phi(Y)}
 =\gamma\inner{Z}{Y},
 \qquad Z,Y\in\mathfrak m.
\end{equation}

For $a\in W$, set
$$
 R_t:=\exp(t(v\wedge a)),\qquad v_t:=R_tv,
$$
and use $R_t^{-1}$ to identify $v_t^\perp$ with $W$.  Equivalently, put
$$
 \widetilde\varphi_a(t):=R_t^{-1}\cdot\varphi(x,v_t)
 \in\Lambda^pW^*.
$$
The fixed orbit-type assumption gives
$\widetilde\varphi_a(t)\in\SO(W)\cdot\phi$, so there is a unique
$\xi_a\in\mathfrak m$ such that
$$
 \widetilde\varphi_a'(0)=\Phi_\phi(\xi_a).
$$
Since
$$
 \widetilde\varphi_a'(0)
 =d\varphi_v(a)-\rho_{\Lambda^p}(v\wedge a)\phi,
$$
where $d\varphi_v(a)$ is taken in the fixed ambient space
$\Lambda^pT_x^*M$, we obtain
\begin{equation}\label{eq:vertical-orbit-derivative}
 d\varphi_v(a)
 =\rho_{\Lambda^p}(v\wedge a)\phi+\Phi_\phi(\xi_a)
 =-v^\flat\wedge\iota_a\phi+\Phi_\phi(\xi_a).
\end{equation}
The resulting tensor $\xi\in W^*\otimes\mathfrak m$ is the
$\mathfrak m$-component of the vertical derivative of the orbit-valued map.

For an algebraic curvature operator $\cA$ on $T_xM$, define
\begin{align}
 J_{\cA}(a,b)
 &:=\cA(v\wedge a,v\wedge b),
 \label{eq:orbit-Jacobi-block}\\
 C_{\cA}(a)
 &:=P_{\mathfrak m}P_{\Lambda^2W}\bigl(\cA(v\wedge a)\bigr),
 \notag
\end{align}
and put
\begin{align}
 \mathscr A_{\cA}(\xi)
 &:=\sum_{a,b}J_{\cA}(e_a,e_b)\inner{\xi_a}{\xi_b},
 \label{eq:orbit-A-form}\\
 \mathscr B_{\cA}(\xi)
 &:=\sum_a\inner{C_{\cA}(e_a)}{\xi_a},
 \notag\\
 \Theta_{\cA}
 &:=\operatorname{tr}_{\mathfrak m}
 \bigl(P_{\mathfrak m}\cA|_{\mathfrak m}\bigr).
 \label{eq:orbit-trace-form}
\end{align}

\begin{proposition}[Curvature identity for a fixed orbit type]
\label{prop:orbit-type-curvature}
Assume that \eqref{eq:orbit-homothety} holds.  Then
\begin{equation}\label{eq:orbit-type-Q-evaluation}
 \inner{\cQ_{\mathscr E_{k,p}}(\cA)T}{T}_{\mathrm{ev},x}
 =\gamma\int_{S_xM}
 \left(
  \mathscr A_{\cA}(\xi)-2\mathscr B_{\cA}(\xi)+\Theta_{\cA}
 \right)
 \,d\sigma_x(v).
\end{equation}
Consequently, for $\cA=\kappa_0^{-1}\Rop^g$,
\begin{equation}\label{eq:orbit-type-integrated-inequality}
 \int_{\SM}
 \left(
  \mathscr A_{\cA}(\xi)-2\mathscr B_{\cA}(\xi)+\Theta_{\cA}
 \right)
 \,d\mu\leq0.
\end{equation}
\end{proposition}

\begin{proof}
The infinitesimal action on the homogeneous tensor has two simple forms.
From \eqref{eq:vertical-orbit-derivative},
\begin{equation}\label{eq:orbit-radial-representation}
 \bigl(\rho(v\wedge a)T\bigr)(v^k)
 =\rho_{\Lambda^p}(v\wedge a)\phi-d\varphi_v(a)
 =-\Phi_\phi(\xi_a).
\end{equation}
If $Z\in\Lambda^2W$, then $Zv=0$, and hence
\begin{equation}\label{eq:orbit-tangential-representation}
 \bigl(\rho(Z)T\bigr)(v^k)
 =\rho_{\Lambda^p}(Z)\phi
 =\Phi_\phi(P_{\mathfrak m}Z).
\end{equation}
Choose orthonormal bases $(v\wedge e_a)$ of $v\wedge W$ and $(Z_\mu)$ of
$\Lambda^2W$.  In the quadratic form defining
$\cQ_{\mathscr E_{k,p}}(\cA)$, the radial--radial block is, by
\eqref{eq:orbit-radial-representation} and
\eqref{eq:orbit-homothety}, equal to
$\gamma\mathscr A_{\cA}(\xi)$.  The two mixed blocks together give
$-2\gamma\mathscr B_{\cA}(\xi)$.  Finally,
\eqref{eq:orbit-tangential-representation} gives the tangential block
$\gamma\Theta_{\cA}$.  Integration over the sphere proves
\eqref{eq:orbit-type-Q-evaluation};
\eqref{eq:orbit-type-integrated-inequality} follows from
Corollary~\ref{cor:exceptional-integral-sign}.
\end{proof}

\subsection{Dimension 8 and the \texorpdfstring{$G_2$}{G2} orbit}
\label{sec:dimension-eight}

Put
\begin{equation*}
  \delta_8^\sharp:=4-2\sqrt3
  =0.535898384862246\ldots.
\end{equation*}
Let $W$ be a Euclidean $7$-space with a $G_2$ three-form $\phi$.  Its
stabilizer decomposition is
\begin{equation*}
  \Lambda^2W=\mathfrak g_2\oplus\mathfrak m,
  \qquad \mathfrak m=\Lambda^2_7W,
  \qquad \dim\mathfrak m=7.
\end{equation*}
The orbit differential is a homothety on $\mathfrak m$, so
Proposition~\ref{prop:orbit-type-curvature} applies.

\begin{lemma}[$G_2$ trace and mixed-block estimates]
\label{lem:G2-block-estimates}
Let $V=\bbR v\oplus W$ and let $\cA$ be an algebraic curvature operator on
$V$ satisfying
$$
  \delta\leq\sec_{\cA}\leq1.
$$
For the blocks in \eqref{eq:orbit-A-form}--\eqref{eq:orbit-trace-form},
\begin{align}
  \mathscr A_{\cA}(\xi)&\geq\delta|\xi|^2,
  \notag\\
  \Theta_{\cA}&\geq7\delta,
  \notag\\
  |C_{\cA}|_{\mathrm{HS}}^2
  &\leq\frac{28}{3}(1-\delta)^2.
  \label{eq:G2-C-upper}
\end{align}
\end{lemma}

\begin{proof}
The first inequality follows from the Jacobi lower bound.

Let $\psi=*_{W}\phi$ and define
$$
  \mathcal T_\phi(\alpha):=*_{W}(\phi\wedge\alpha).
$$
The standard $G_2$ identities give eigenvalue $2$ on $\Lambda^2_7$ and
$-1$ on $\Lambda^2_{14}$, hence
$$
  P_{\mathfrak m}=\frac13(\Id+\mathcal T_\phi).
$$
The operator $\mathcal T_\phi$ is the curvature-operator realization of the
$4$-form $\psi$.  Since an algebraic curvature operator is orthogonal to
$4$-forms by the first Bianchi identity,
$$
  \Theta_{\cA}
  =\tr_{\mathfrak m}(P_{\mathfrak m}\cA|_{\mathfrak m})
  =\frac13\tr_{\Lambda^2W}(\cA)
  \geq\frac13\binom72\delta=7\delta.
$$

Choose
$$
  Z_0=\frac1{\sqrt3}(e_{23}+e_{45}+e_{67})\in\Lambda^2_7W.
$$
The $G_2$-action on the unit sphere of $\Lambda^2_7W$ is transitive, and
Schur's lemma gives the tight-frame identity
$$
  P_{\mathfrak m}
  =7\int_{G_2}(gZ_0)\otimes(gZ_0)\,dg.
$$
Consequently
$$
  |C_{\cA}|_{\mathrm{HS}}^2
  =7\int_{G_2}|C_{\cA}^*(gZ_0)|^2\,dg.
$$
Every $gZ_0$ is a normalized sum of three mutually orthogonal unit simple
two-forms.  The constant-curvature operator has no radial--tangential
block, so Lemma~\ref{lem:BBK} gives, for every unit $a\in W$,
$$
  |\cA(v\wedge a,gZ_0)|
  \leq\frac1{\sqrt3}\,3\,\frac23(1-\delta)
  =\frac2{\sqrt3}(1-\delta).
$$
Thus $|C_{\cA}^*(gZ_0)|^2\leq\frac43(1-\delta)^2$, which proves
\eqref{eq:G2-C-upper}.
\end{proof}

\begin{proposition}
\label{prop:G2-eight-vanishing}
Let $(M^8,g)$ be closed and $\delta$-pinched for some
$\delta>\delta_8^\sharp$.  Then there is no odd flow-invariant
$G_2$-structure
$$
  \phi\in C^\infty(\SM,\Lambda^3\cN).
$$
\end{proposition}

\begin{proof}
Apply Lemma~\ref{lem:finite-fourier-homogenization} to the full finite odd
Fourier expansion of $\phi$, and let $\xi$ be the tensor defined by
\eqref{eq:vertical-orbit-derivative}.  Proposition~\ref{prop:orbit-type-curvature} gives the integrated
orbit-type curvature inequality.  Lemma~\ref{lem:G2-block-estimates} and
Cauchy--Schwarz give pointwise
\begin{align*}
 \mathscr A_{\cA}(\xi)-2\mathscr B_{\cA}(\xi)+\Theta_{\cA}
 &\geq
 \delta|\xi|^2-2|C_{\cA}|_{\mathrm{HS}}|\xi|+7\delta\\
 &\geq
 7\delta-\frac{|C_{\cA}|_{\mathrm{HS}}^2}{\delta}\\
 &\geq
 7\delta-\frac{28(1-\delta)^2}{3\delta}\\
 &=\frac7{3\delta}\bigl(-\delta^2+8\delta-4\bigr)>0,
\end{align*}
where the last inequality is equivalent to
$\delta>4-2\sqrt3$.  This contradicts
\eqref{eq:orbit-type-integrated-inequality}.
\end{proof}

\begin{lemma}
\label{lem:projector-eight}
Let $(M^8,g)$ be closed and $5/13$-pinched.  There is no nontrivial even
flow-invariant orthogonal projector in
$C^\infty(\SM,\Sym^2\cN)$ of rank at most three.
\end{lemma}

\begin{proof}
The high-degree proof of Proposition~\ref{prop:sym2-high-degree} remains
valid when $n=8$: writing $k=6+b$ in \eqref{eq:sym2-Psi} gives
$$
  \Psi_{8,6+b}
  =450b^3+10212b^2+63756b+70128>0.
$$
For completeness, Weyl's dimension formula at $n=8$ gives
$$
\begin{array}{c|rrrrrr}
 \lambda&(6)&(5,1)&(4,2)&(4)&(3,1)&(2)\\ \hline
 \dim\mathbb V^0_\lambda&1386&3696&4312&294&567&35.
\end{array}
$$
Their sum is
$$
 10290=294\cdot35
 =\dim\Harm_4(\bbR^8)\,\dim\Sym^2_0(\bbR^8),
$$
so the full orthogonal Pieri decomposition has neither an additional
low-rank summand nor a multiplicity.  Likewise,
$1386+3696+294=5376=8\dim\Harm_5(\bbR^8)$; the same highest-weight
contractions are nonzero, and the explicit embeddings
\eqref{eq:SH-embedding}--\eqref{eq:Sq-embedding} remain valid.  Hence
\eqref{eq:quartic-decomposition} holds for $O(8)$.  The three gaps in
\eqref{eq:quartic-gaps} become
$$
  \frac{77}{18},\qquad \frac{94}{15},\qquad \frac{193}{54}.
$$
Thus degrees at least $4$ are impossible.  Proposition
\ref{prop:quadratic-projector}, which is valid in every even dimension,
excludes the remaining degree-two projector of rank $1\leq r\leq3$.
\end{proof}

\begin{corollary}
\label{cor:dimension-eight-threshold}
If $(M^8,g)$ is closed, oriented, negatively curved, and strictly
$\delta_8^\sharp$-pinched, then its oriented frame flow is ergodic.
\end{corollary}

\begin{proof}
If the frame flow were not ergodic, Theorem~\ref{thm:CLMS-input} would give,
after a finite cover, either an odd invariant $G_2$-structure or an even
invariant projector of rank at most three.  Proposition
\ref{prop:G2-eight-vanishing} excludes the first.  Since
$4-2\sqrt3>5/13$, Lemma~\ref{lem:projector-eight} excludes the second.
\end{proof}

\subsection{Dimension 134 and the \texorpdfstring{$E_7$}{E7} Cartan orbit}
\label{sec:dimension-134}

Assume that the dynamical reduction supplies a nonzero flow-invariant Lie
bracket
\begin{equation*}
 \varphi\in C^\infty(\SM,\Lambda^3\cN),
 \qquad X\varphi=0.
\end{equation*}
At every $(x,v)$, the three-form $\varphi(x,v)$ is equivalent to the Cartan
three-form
\begin{equation*}
 \varphi_0(a,b,c)=\inner{[a,b]}c
\end{equation*}
on the compact Lie algebra $\mathfrak e_7$.  Its stabilizer in
$\mathrm O(133)$ is $E_7/\{\pm1\}\subset\SO(133)$; see
\cite[Theorem~3.8 and Lemma~3.12]{CLMS24}.

Fix $(x,v)\in\SM$, put $W=v^\perp$, and abbreviate
$\phi:=\varphi(x,v)$.  The bracket identifies a copy
$\mathfrak h\simeq\mathfrak e_7$ inside $\Lambda^2W$.  After
complexification, \cite[Proposition~5.1]{Reeder97} gives
$$
 \Lambda^2(\mathfrak e_7\otimes\bbC)
 \simeq(\mathfrak e_7\otimes\bbC)\oplus U_2,
$$
where $U_2$ is irreducible.  Since
$\dim U_2=\binom{133}{2}-133=8645$, the compact real form has the
irreducible orthogonal decomposition
\begin{equation*}
 \Lambda^2W=\mathfrak h\oplus\mathfrak m,
 \qquad
 \mathfrak h\simeq\mathfrak e_7,
 \qquad
 \dim\mathfrak m=8645.
\end{equation*}
Hence the orbit differential \eqref{eq:orbit-differential} is a homothety
and the orbit-type curvature identity applies.

\begin{lemma}
\label{lem:E7-tight-frame}
Let $\cA$ be an algebraic curvature operator on $V=\bbR v\oplus W$ satisfying
$$
 \delta\leq\sec_{\cA}\leq1.
$$
For the quantities in
\eqref{eq:orbit-Jacobi-block}--\eqref{eq:orbit-trace-form}, one has
\begin{align}
 \mathscr A_{\cA}(\xi)&\geq\delta|\xi|^2,
 \notag\\
 \Theta_{\cA}&\geq8645\,\delta,
 \notag\\
 |C_{\cA}|_{\mathrm{HS}}^2
 &\leq\frac{4\cdot8645}{9}(1-\delta)^2.
 \label{eq:E7-C-upper}
\end{align}
\end{lemma}

\begin{proof}
The first inequality follows because the Jacobi operator satisfies
$J_{\cA}\succeq\delta\Id_W$.

Identify $W$ with the compact Lie algebra $\mathfrak e_7$ using $\phi$.
Choose orthonormal commuting vectors $p,q$ in a Cartan subalgebra and put
$z=p\wedge q$.  The copy of $\mathfrak h$ in $\Lambda^2W$ is generated by
the two-forms $\iota_u\phi$, $u\in W$.  Since $[p,q]=0$,
$$
 \inner{p\wedge q}{\iota_u\phi}
 =\phi(u,p,q)=\inner{[u,p]}q
 =\inner{u}{[p,q]}=0,
$$
so $z\in\mathfrak m$.  It is a unit decomposable two-form.

Let $H=E_7/\{\pm1\}$ and normalize its Haar measure to have mass one.
Irreducibility of $\mathfrak m$ and Schur's lemma give
\begin{equation*}
 P_{\mathfrak m}
 =8645\int_H(gz)\otimes(gz)\,dg.
\end{equation*}
Taking traces fixes the coefficient.  Therefore
$$
 \Theta_{\cA}
 =8645\int_H\cA(gz,gz)\,dg
 \geq8645\,\delta.
$$
Similarly,
\begin{equation}\label{eq:E7-C-tight-frame}
 |C_{\cA}|_{\mathrm{HS}}^2
 =8645\int_H|C_{\cA}^*(gz)|^2\,dg.
\end{equation}
Write $gz=(gp)\wedge(gq)$.  For a unit vector $x\in W$,
$$
 \inner{C_{\cA}^*(gz)}{x}=\cA(v,x,gp,gq).
$$
The constant-curvature operator has zero pairing between the radial simple
form $v\wedge x$ and the tangential simple form $(gp)\wedge(gq)$.  Hence
Lemma~\ref{lem:BBK} gives
$$
 |\cA(v,x,gp,gq)|\leq\frac23(1-\delta).
$$
Taking the supremum over $x$ and using
\eqref{eq:E7-C-tight-frame} proves \eqref{eq:E7-C-upper}.
\end{proof}

\begin{proposition}
\label{prop:E7-vanishing}
Let $(M^{134},g)$ be closed and $\delta$-pinched for some
$\delta>2/5$.  Then there is no nonzero flow-invariant Lie bracket
$$
 \varphi\in C^\infty(\SM,\Lambda^3\cN)
$$
whose values are equivalent to the Cartan three-form of the compact Lie
algebra $\mathfrak e_7$.
\end{proposition}

\begin{proof}
Apply Lemma~\ref{lem:finite-fourier-homogenization} to $\varphi$ and let $\xi$
be the tensor defined by \eqref{eq:vertical-orbit-derivative}.  By
Lemma~\ref{lem:E7-tight-frame} and
Cauchy--Schwarz,
\begin{align}
 \mathscr A_{\cA}(\xi)-2\mathscr B_{\cA}(\xi)+\Theta_{\cA}
 &\geq
 \delta|\xi|^2-2|C_{\cA}|_{\mathrm{HS}}|\xi|+8645\,\delta
 \notag\\
 &\geq
 8645\,\delta-\frac{|C_{\cA}|_{\mathrm{HS}}^2}{\delta}
 \notag\\
 &\geq
 8645\left(
  \delta-\frac{4(1-\delta)^2}{9\delta}
 \right)
 \notag\\
 &=8645\,
 \frac{(5\delta-2)(\delta+2)}{9\delta}>0.
 \label{eq:E7-final-positive}
\end{align}
This pointwise strict positivity contradicts
\eqref{eq:orbit-type-integrated-inequality}.
\end{proof}

\begin{remark}[Endpoint of the orbit-type curvature estimate]
\label{rem:E7-endpoint}
At $\delta=2/5$ the final $E_7$ lower bound in
\eqref{eq:E7-final-positive} vanishes.  Likewise, the $G_2$ lower bound
vanishes at $\delta=4-2\sqrt3$. 
\end{remark}

\subsection{Dimension 7: a cubic refinement}
\label{sec:dimension-seven}

Put
\begin{equation*}
  \delta_7^\sharp
  :=\frac{385+36\sqrt6}{927+36\sqrt6}
  =0.466105390810889\ldots.
\end{equation*}
Throughout this subsection, norms are the standard fibrewise or global
$L^2$ norms.  For a degree-two vector harmonic, write
$\mathcal W_{\cB}:=c_{7,2}^{-1}W_{\cB}$; this removes the harmless algebraic
normalization in \eqref{eq:tangent-twist-form} and is the normalization
appearing directly in the localized Pestov identity.

\begin{proposition}[The $7$-dimensional cubic refinement]
\label{prop:dimension-seven-vanishing}
Let $(M^7,g)$ be closed and negatively curved.  If $g$ is strictly
$\delta_7^\sharp$-pinched, then there is no odd flow-invariant orthogonal
complex structure
$$
  f\in C^\infty(\SM,\Lambda^2\cN).
$$
Consequently the oriented frame flow is ergodic.
\end{proposition}

\begin{proof}
Choose a pinching presentation with
$\delta>\delta_7^\sharp$ and
$\cA=\delta I+(1-\delta)\cE$, where
$0\leq\sec_{\cE}\leq1$.  Let $k$ be the highest odd Fourier degree of
$f$.  The estimates of \cite[Proposition~4.4 and Lemma~4.11]{CLMS24}
already exclude every $k\geq5$ under $\delta\geq2/5$: at
$(n,p,k)=(7,2,5)$ the two relevant gaps are
$$
  \frac{93}{5}-8\sqrt3>0,
  \qquad
  \frac{2807}{108}-8\sqrt3-\frac{26}{9}\sqrt6>0,
$$
and the corresponding thresholds decrease with $k$.

It remains to exclude $k=3$.  Put $u=f_3$ and $g=\iota_vu$.  We use the
localized Pestov calculation of \cite[Proposition~4.4]{CLMS24}, but retain
the degree-two vector term.  At $(n,p,k)=(7,2,3)$ the lowering coefficient
is
$$
 \frac{(n+k-2)(n+2k-4)}{n+k-3}=\frac{72}{7}.
$$
After moving the curvature terms to the left, the Jacobi term contributes
$24\delta\|u\|^2$, the constant-curvature part of the $\Lambda^2$ twist
subtracts
$$
 (1+\delta)\|u\|^2+\frac92(1+\delta)\|g\|^2,
$$
and the centered twist is bounded by $16(1-\delta)\|u\|^2$.  Discarding
the nonnegative $Z(u)$ term and using
$$
 24\delta-(1+\delta)-16(1-\delta)=39\delta-17
$$
gives
\begin{equation}\label{eq:seven-cubic-master}
 0\geq(39\delta-17)\|u\|^2
 -\frac92(1+\delta)\|g\|^2
 +\frac{72}{7}\|X_-u\|^2.
\end{equation}
Since $X_+u=0$ and tautological contraction commutes with $X$,
$$
  \|X_-u\|^2\geq\|Xg\|^2\geq\|X_+g\|^2.
$$
For a degree-two vector harmonic, the localized Pestov identity gives
$$
  \mathcal W_{\cA}(g)\leq\frac{22}{3}\|X_+g\|^2.
$$
Lemma
\ref{lem:seven-degree-two-estimate} therefore gives
$$
  \|X_+g\|^2
  \geq\frac3{22}
  \bigl(13\delta-2\sqrt6(1-\delta)\bigr)\|g\|^2.
$$
Substitution in \eqref{eq:seven-cubic-master} yields
\begin{equation}\label{eq:seven-cubic-final-inequality}
  0\geq B_7(\delta)\|u\|^2+C_7(\delta)\|g\|^2,
\end{equation}
where
$$
  B_7(\delta)=39\delta-17,
$$
and
$$
  C_7(\delta)
  =\frac{108}{77}
    \bigl(13\delta-2\sqrt6(1-\delta)\bigr)
   -\frac92(1+\delta).
$$
Now $B_7(\delta)>0$.  If $C_7(\delta)\geq0$,
\eqref{eq:seven-cubic-final-inequality} is already impossible for $u\neq0$.
If $C_7(\delta)<0$, Lemma~\ref{lem:seven-contraction} gives
$$
  B_7(\delta)\|u\|^2+C_7(\delta)\|g\|^2
  \geq
  \left(B_7(\delta)+\frac23C_7(\delta)\right)\|u\|^2.
$$
A direct simplification gives
\begin{equation*}
  B_7(\delta)+\frac23C_7(\delta)
  =\frac4{77}
   \left((927+36\sqrt6)\delta-(385+36\sqrt6)\right)>0.
\end{equation*}
Thus the cubic coefficient vanishes.

The only remaining possibility is degree one.  Then
$f(v)=\iota_v\phi$ for a three-form $\phi$ on $M$, and the condition that
$f(v)$ be an orthogonal complex structure on $v^\perp$ makes $\phi$ a
$G_2$-structure.  The equation $Xf=0$ makes this structure nearly parallel;
such a metric is Einstein with nonnegative scalar curvature, contradicting
negative sectional curvature.  This final degree-one argument is the one
recorded in \cite[proof of Theorem~4.1]{CLMS24}.  Theorem
\ref{thm:CLMS-input} now implies ergodicity.
\end{proof}

\section{Proofs of the main theorems}\label{sec:mainproof}

\begin{proof}[Proof of Theorem \ref{thm:stable-main}]
Assume first that we are in one of the dimensions of
Theorem \ref{thm:stable-main}\textup{(i)} and that the oriented frame flow
is not ergodic.  By Theorem \ref{thm:CLMS-input}, after passing to a finite
Riemannian cover we obtain an odd invariant unit normal vector field $f$.
Strict quarter-pinching pulls back to the cover. Choose a pinching
presentation with $\delta>1/4$ and scale $\kappa_0$, and put
$\cA=\kappa_0^{-1}\Rop^g$.

By Proposition \ref{prop:degree-reduction},
$$
  f=f_1+f_3,
  \qquad f_3\neq0.
$$
Construct $T\in C^\infty(M,\cC)$ by \eqref{eq:Tdef}.  Lemmas
\ref{lem:young-symmetry} and \ref{lem:killing-equation} give
$$
  T\in C^\infty(M,\cC),
  \qquad
  \cD T=0.
$$

If $n=4$, Proposition \ref{prop:curvature-positive} makes
$\cQ_{\Hook}(\cA)$ pointwise positive definite. Compactness gives $c>0$ such
that
$$
  \inner{\cQ_{\Hook}(\cA)T}{T}\geq c|T|^2.
$$
Together with \eqref{eq:integral-negative-Q}, this implies $T=0$.

If $n=6$, let $K$ be the cubic coefficient corresponding to $f_3$.
Lemma \ref{lem:six-dimensional-cubic} and
\eqref{eq:tangent-twist-pestov-normalized} imply $K=0$, contradicting
$f_3\neq0$.

If $n\equiv2\pmod4$ and $n\geq10$, Proposition
\ref{prop:stable-cubic} gives $T=0$.  In the first and third cases,
\eqref{eq:diagonal-T} then gives $f=0$; in the second case we already have a
contradiction to $f_3\neq0$.  This proves part \textup{(i)}.

Now suppose that $n\equiv0\pmod4$, $n\geq12$, and that $g$ is
$5/13$-pinched.  If the frame flow were not ergodic, Theorem
\ref{thm:CLMS-input} would give, on a finite Riemannian cover, a nontrivial
even invariant orthogonal projector
$$
  P\in C^\infty(\SM,\Sym^2\cN).
$$
Choose a normalized pinching presentation with
$\frac5{13}\leq\sec_{\cA}\leq1$.  Proposition
\ref{prop:symmetric-projector-vanishing} excludes precisely such a
projector.  This contradiction proves part \textup{(ii)}.
\end{proof}

\begin{proof}[Proof of Theorem \ref{thm:exceptional-main}]
Part \textup{(i)} is Proposition~\ref{prop:dimension-seven-vanishing}, and
part \textup{(ii)} is Corollary~\ref{cor:dimension-eight-threshold}.

For part \textup{(iii)}, let $n=134$ and assume that $g$ is strictly
$2/5$-pinched.  If the frame flow were not ergodic,
Theorem~\ref{thm:CLMS-input} would give on a finite Riemannian cover either
an odd invariant unit normal vector field or an invariant $E_7$
Lie-bracket three-form.  In the first case strict $2/5$-pinching implies
strict quarter-pinching, and Propositions \ref{prop:degree-reduction} and
\ref{prop:stable-cubic} eliminate the vector field.  In the second case
Proposition~\ref{prop:E7-vanishing} eliminates the Lie-bracket three-form.
This proves part \textup{(iii)}.
\end{proof}



\appendix

The next two appendices record the two short algebraic verifications used above, so
that the paper remains self-contained.

\section{Jucys--Murphy and conformal weights}\label{sec:weitzenbock}\label{app:JM}

This appendix proves the differential half of the argument. We work first in a
Euclidean vector space $V$ of arbitrary dimension at least $4$ and put
$$
  E:=\Hook V^*.
$$



\subsection{The conformal-weight operator}

For $a\in V^*$ and $u\in E$, define
\begin{equation*}
  B(a\otimes u)
  :=\sum_i e^i\otimes\rho(e_i\wedge a)u,
\end{equation*}
where $(e_i)$ is an orthonormal basis and $\rho$ is the infinitesimal
$\mathfrak{so}(V)$-action on $E$. This definition is independent of the basis.

Realise $V^*\otimes E$ inside $V^{*\otimes5}$. Slot $0$ is the first factor and
slots $1,2,3,4$ are the $4$ tensor slots of $E$. Let $\tau_{0r}$ exchange slots
$0$ and $r$, and let $\operatorname{tr}_{0r}$ denote contraction of these
slots.

\begin{lemma}\label{lem:Bformula}
On $V^*\otimes E$,
\begin{equation}\label{eq:Bformula}
  B=\sum_{r=1}^4
  \left(
    \tau_{0r}-\operatorname{tr}_{0r}^*\operatorname{tr}_{0r}
  \right).
\end{equation}
\end{lemma}

\begin{proof}
It is enough to compute on a pure covariant tensor
$a\otimes u_1\otimes\cdots\otimes u_4$. The infinitesimal action of
$e_i\wedge a$ on the $r$-th covariant factor is
$$
  u_r\longmapsto
  \inner{e_i}{u_r}a-\inner{a}{u_r}e_i.
$$
After summing over $i$, the first term exchanges $a$ and $u_r$, while the
second contracts them and inserts the metric tensor. This gives
$\tau_{0r}-\operatorname{tr}_{0r}^*\operatorname{tr}_{0r}$ in the $r$-th
slot. Summing over the $4$ slots proves the formula. Restriction to the
Young-symmetry subspace $E$ is legitimate because all terms are
$\mathrm O(V)$-equivariant.
\end{proof}

\subsection{The Jucys--Murphy eigenvalue calculation}

Over $\bbC$, Pieri's rule gives
\begin{equation}\label{eq:pieri}
  V_\bbC^*\otimes\Hook V_\bbC^*
  \simeq
  \mathbb S_{(4,1)}V_\bbC^*
  \oplus
  \mathbb S_{(3,2)}V_\bbC^*
  \oplus
  \mathbb S_{(3,1,1)}V_\bbC^*.
\end{equation}
The principal symbol of $\cD$ is the projection onto
$\mathbb S_{(4,1)}V_\bbC^*$. Hence
\begin{equation}\label{eq:kernel-symbol}
  \ker\sigma(\cD)
  =
  \mathbb S_{(3,2)}V_\bbC^*
  \oplus
  \mathbb S_{(3,1,1)}V_\bbC^*.
\end{equation}

After relabelling the distinguished tensor slot $0$ as $5$, set
$$
  J_5:=\sum_{r=1}^4\tau_{0r}.
$$
Under the place-permutation action of $\bbC[S_5]$ on
$V_\bbC^{*\otimes5}$, this is the image of the fifth Jucys--Murphy element
$$
  (1\,5)+(2\,5)+(3\,5)+(4\,5).
$$
On the Pieri summand obtained by adding one box to $(3,1)$,
$J_5$ acts by the content
$\operatorname{column}-\operatorname{row}$ of the added box;
see \cite[\S\S3.2--3.3]{CeccheriniScarabottiTolli10}.
Thus its eigenvalues on the three summands in \eqref{eq:pieri} are
\begin{equation}\label{eq:contents}
  3,\qquad 0,\qquad -2,
\end{equation}
respectively; see also Subsection~\ref{app:JM-verification}.

\begin{proposition}\label{prop:symbol-sign}
If $U\in V^*\otimes E$ satisfies $\sigma(\cD)U=0$, then
$$
  \inner{BU}{U}\leq0.
$$
\end{proposition}

\begin{proof}
By \eqref{eq:kernel-symbol} and \eqref{eq:contents},
$$
  \inner{J_5U}{U}\leq0.
$$
On the other hand, Lemma \ref{lem:Bformula} gives
$$
  \inner{BU}{U}
  =\inner{J_5U}{U}
   -\sum_{r=1}^4\norm{\operatorname{tr}_{0r}U}^2
  \leq\inner{J_5U}{U}.
$$
The conclusion follows.
\end{proof}

\subsection{Integrated Weitzenb\"ock formula}

Let $q_E(R^g)$ denote the Weitzenb\"ock curvature endomorphism on the bundle
associated with $E$. For an equivariant endomorphism
$B\in\End(V^*\otimes E)$, let
$$
  \widetilde B(a\otimes b\otimes u)
  :=(a\mathbin{\lrcorner}\otimes\Id)B(b\otimes u),
$$ where \(a\mathbin{\lrcorner}\) denotes contraction of the
first \(V^*\)-factor with \(a^\sharp\), using the Euclidean metric.
We use $(\nabla^2T)(a,b):=\nabla_a\nabla_bT-\nabla_{\nabla_a b}T$.
The universal conformal-weight identity is
\begin{equation*}
  \widetilde B(\nabla^2T)=q_E(R^g)T;
\end{equation*}
see \cite[Lemma~3.1]{CLMSW25}. Since $B$ is parallel and self-adjoint,
integration by parts on a closed manifold yields
\begin{equation}\label{eq:integrated-W}
  \int_M\inner{q_E(R^g)T}{T}\,d\vol_g
  =-\int_M\inner{B\nabla T}{\nabla T}\,d\vol_g.
\end{equation}

\begin{corollary}\label{cor:integral-sign}
If $T\in C^\infty(M,\cC)$ satisfies the generalized Killing equation
\eqref{eq:generalized-killing}, then
\begin{equation*}
  \int_M\inner{q_E(R^g)T}{T}\,d\vol_g\geq0.
\end{equation*}
\end{corollary}


\begin{proof}
The equation $\cD T=0$ says pointwise that
$\nabla T\in\ker\sigma(\cD)$. Apply Proposition \ref{prop:symbol-sign} in
\eqref{eq:integrated-W}, and then use
\eqref{eq:q-Q-exact-convention}.
\end{proof}


Equivalently,
\begin{equation*}
  -\int_M\inner{\cQ_E(\Rop^g)T}{T}\,d\vol_g\geq0.
\end{equation*}
If $g$ is $\delta$-pinched at scale $\kappa_0$ and $\cA$ is the normalized
operator from \eqref{eq:normalized-curvature}, this becomes
\begin{equation}\label{eq:integral-negative-Q}
  -\kappa_0\int_M\inner{\cQ_E(\cA)T}{T}\,d\vol_g\geq0.
\end{equation}

\subsection{Verification of the Jucys--Murphy eigenvalues}\label{app:JM-verification}

For clarity, we recall why the eigenvalues in \eqref{eq:contents} are
$3,0,-2$. Let $[\lambda]$ be the irreducible symmetric-group module associated
with a partition $\lambda$. The restriction/induction branching rule says
that inducing $[(3,1)]$ from $S_4$ to $S_5$ gives the direct sum of the modules
obtained by adding one admissible box:
$$
  [(4,1)]\oplus[(3,2)]\oplus[(3,1,1)].
$$
The Jucys--Murphy element
$$
  J_5=(1\,5)+(2\,5)+(3\,5)+(4\,5)
$$
acts on a Young basis by the content of the box containing $5$. The added boxes
have coordinates
$$
  (1,4),\qquad(2,2),\qquad(3,1),
$$
and hence contents
$$
  4-1=3,
  \qquad
  2-2=0,
  \qquad
  1-3=-2.
$$
Schur--Weyl duality identifies the same scalars on the three $GL(V)$ summands
in \eqref{eq:pieri}.

\section{The first-three-slot contraction}\label{app:first-three-slot}

For completeness we record the word calculation used in
Lemma \ref{lem:first-three-slot-action}. For $L\in\{\mathsf D,\mathsf O\}$, write
$L_{ab}:=L_{e_a\wedge e_b}$. For
$L_1,L_2,L_3\in\{\mathsf D,\mathsf O\}$, let
$$
 \mathsf W^{L_1L_2L_3}_{ab}
 :=\frac12\operatorname{pr}_{\cU}\sum_{i,j}
 \bigl((L_1)_{ai}(L_2)_{ij}(L_3)_{jb}
      -(L_1)_{bi}(L_2)_{ij}(L_3)_{ja}\bigr)
 \operatorname{pr}_{\cU},
$$
and define the two-letter expressions analogously. We spell out two
representative contractions. Let $\mathsf E_{ab}$ be the matrix unit on the
fourth covariant slot, so that
$\mathsf O_{ab}=\mathsf E_{ba}-\mathsf E_{ab}$. A direct matrix-unit
contraction gives
\begin{align*}
 \frac12\sum_i
 \bigl(\mathsf O_{ai}\mathsf O_{ib}
      -\mathsf O_{bi}\mathsf O_{ia}\bigr)
 &=-\frac{n-2}{2}\mathsf O_{ab},\\
 \frac12\sum_{i,j}
 \bigl(\mathsf O_{ai}\mathsf O_{ij}\mathsf O_{jb}
      -\mathsf O_{bi}\mathsf O_{ij}\mathsf O_{ja}\bigr)
 &=\frac{n^2-3n+4}{2}\mathsf O_{ab}.
\end{align*}
Since
$\operatorname{pr}_{\cU}\mathsf O_{ab}\operatorname{pr}_{\cU}
=\rho_{ab}-\mathsf M_{ab}$, these are precisely the
$\mathsf{OO}$ and $\mathsf{OOO}$ rows below. For words involving
$\mathsf D$, one inserts \eqref{eq:tracefree-hook-projection} after each word and
uses $h_{(abcd)}=0$ together with the vanishing traces. For example,
$$
 \mathsf W^{\mathsf{DD}}_{ab}=-\frac{n-2}{2}\mathsf M_{ab},
 \qquad
 \mathsf W^{\mathsf{DDO}}_{ab}
 =\frac{n-4}{2}(\mathsf M_{ab}-\rho_{ab}),
 \qquad
 \mathsf W^{\mathsf{DOD}}_{ab}=-\mathsf M_{ab}.
$$
The remaining words are obtained by the same contraction. Thus
\eqref{eq:tracefree-hook-projection} gives the following coefficients in the
basis $(\mathsf M_{ab},\rho_{ab})$:
$$
\begin{array}{c|cc@{\qquad}c|cc}
\text{word}&\mathsf M&\rho&\text{word}&\mathsf M&\rho\\ \hline
\mathsf{DD}&-(n-2)/2&0&\mathsf{DDD}&(n^2+n+12)/2&0\\
\mathsf{DO},\mathsf{OD}&0&0&\mathsf{DDO},\mathsf{ODD}&(n-4)/2&-(n-4)/2\\
\mathsf{OO}&(n-2)/2&-(n-2)/2&\mathsf{DOD}&-1&0\\
&&&\mathsf{DOO},\mathsf{OOD}&n/2&-(n-2)/2\\
&&&\mathsf{ODO}&1&-1\\
&&&\mathsf{OOO}&-(n^2-3n+4)/2&(n^2-3n+4)/2.
\end{array}
$$
The two-letter row sums to \eqref{eq:alt-B2}; the eight three-letter words
sum to \eqref{eq:alt-B3}. Each entry uses only the first Bianchi/Young
relation, $\delta_{ii}=n$, the vanishing traces of a tensor in the totally trace-free $(3,1)$-component, and the two
eigenvalues of $CC^*$ displayed before
\eqref{eq:tracefree-hook-projection}.

\section{The quartic admissible module and eigenvalue table}\label{app:quartic-module}

Let $\mathscr N^{(2)}_4$ denote the space of quartic harmonic
$\Sym^2V$-valued polynomials satisfying
\eqref{eq:sym2-normal-conditions}. In the range $n\geq12$ used below, the
orthogonal Pieri rule gives the multiplicity-free decomposition
\begin{equation}\label{eq:quartic-full-pieri}
\begin{split}
 \Harm_4(V)\otimes\Sym^2_0V
 \simeq{}&
 \mathbb V^0_{(6)}V^*\oplus\mathbb V^0_{(5,1)}V^*
 \oplus\mathbb V^0_{(4,2)}V^*\\
 &\oplus\mathbb V^0_{(4)}V^*
 \oplus\mathbb V^0_{(3,1)}V^*
 \oplus\mathbb V^0_{(2)}V^*.
\end{split}
\end{equation}
Let $\pi_5(S)$ be the degree-five harmonic component of
$x\mapsto S(x)x$. Orthogonal Pieri also gives
$$
 \Harm_5(V)\otimes V
 \simeq
 \mathbb V^0_{(6)}V^*\oplus\mathbb V^0_{(5,1)}V^*
 \oplus\mathbb V^0_{(4)}V^*.
$$
The map $\pi_5$ is equivariant, and evaluation on one highest-weight vector
shows that it is nonzero on each of the three matching summands in
\eqref{eq:quartic-full-pieri}. Multiplicity one therefore gives
$$
 \ker\pi_5
 =\mathbb V^0_{(4,2)}V^*
 \oplus\mathbb V^0_{(3,1)}V^*
 \oplus\mathbb V^0_{(2)}V^*.
$$
The condition that $S(v)v$ have degree three implies
$\mathscr N^{(2)}_4\subseteq\ker\pi_5$. We now verify that all three summands
in this kernel satisfy the remaining normal condition, which will prove
\begin{equation}\label{eq:quartic-decomposition}
 \mathscr N^{(2)}_4
 \simeq
 \mathbb V^0_{(4,2)}V^*
 \oplus\mathbb V^0_{(3,1)}V^*
 \oplus\Sym^2_0V^*.
\end{equation}
We record concrete models for the last two summands.  Let $H$ be a harmonic
homogeneous vector polynomial of degree three satisfying
\begin{equation*}
  \inner{H(x)}x=0,
  \qquad
  \operatorname{div}H=0.
\end{equation*}
Then
\begin{equation}\label{eq:SH-embedding}
 S_H(x)=\frac{n+4}{n+2}
 \bigl(x\otimes H+H\otimes x\bigr)
 -\frac{|x|^2}{n+2}
 \bigl(\nabla H+\nabla H^\top\bigr).
\end{equation}
If $q\in\Harm_2(V)$, put
\begin{align}
 S_q(x)=\frac1{n(n+1)}\bigl[&
 (n+4)(n+2)q\,x\otimes x-(n+2)|x|^2q\,\Id_V\notag\\
 &-(n+2)|x|^2
 (x\otimes\nabla q+\nabla q\otimes x)
 +|x|^4\nabla^2q\bigr].
 \label{eq:Sq-embedding}
\end{align}
Both expressions are harmonic, trace-free, and satisfy the normal
conditions.  Moreover,
$$
  S_H(x)x=|x|^2H(x),
  \qquad
  S_q(v)v=q(v)v-\frac1n\nabla^Sq(v)
  \quad(|v|=1).
$$
The contraction $S\mapsto S(v)v$ vanishes on the $(4,2)$ summand, whereas
the displayed identities show that $S_H$ and $S_q$ define nonzero equivariant
maps from the $(3,1)$ and $(2)$ summands, respectively. Multiplicity one then
identifies these three models with all of $\mathscr N^{(2)}_4$.

\begin{lemma}[Quartic eigenvalue table]
\label{lem:quartic-table}
The values of the $6$ invariant quadratic forms on the three summands of
\eqref{eq:quartic-decomposition} are given in
Table~\eqref{eq:quartic-table}.  The first row is normalized by
$\int|S|^2=1$, the second by $\int|H|^2=1$, and the third by
$\int q^2=1$.
\begin{equation}\label{eq:quartic-table}
\begin{array}{c|cccccc}
 &\mathscr W_I^{(2)}&A_0&A_C&B&Z_b&Z_C\\ \hline
(4,2)
&4n+6
&1
&n-1
&0
&1
&2n+3\\[2mm]
(3,1)
&\dfrac{2(n+4)(3n+4)}{n+2}
&n+\dfrac4{n+2}
&\dfrac{4(n-1)}{n+2}
&1
&\dfrac{3n^2+12n+16}{n+2}
&\dfrac{n^2+10n+12}{n+2}\\[4mm]
(2)
&\dfrac{2(n+2)^2(n+4)}{n(n+1)}
&n+1+\dfrac{n+4}{n(n+1)}
&\dfrac{2(n-2)}{n(n+1)}
&\dfrac2n
&\dfrac{(n+2)^2(n+4)}{n(n+1)}
&\dfrac{3(n+2)^2}{n(n+1)}.
\end{array}
\end{equation}
\end{lemma}

\begin{proof}
We give the calculation, including the derivative energies.  For $k=4$,
\eqref{eq:sym2-WI} reads
\begin{equation}\label{eq:quartic-WI-common}
 \mathscr W_I^{(2)}=(4n+6)N-2(n+4)N_1.
\end{equation}
If
$$
 Z_{00}:=\int|z_{00}|^2,
$$
then the spherical harmonic energy splits as
\begin{equation}\label{eq:quartic-energy-split}
 4(n+2)N=Z_{00}+2Z_b+2Z_C.
\end{equation}

On the $(4,2)$ summand one has $a=b=0$ and $C=S$.  Differentiating
$S(v)v=0$ tangentially gives $z_{00}=0$ and $Z_b=N$.  Equations
\eqref{eq:quartic-WI-common} and \eqref{eq:quartic-energy-split} then give
the first row.

For $S_H$, one has on the unit sphere
$$
 a=0,
 \qquad b=H,
 \qquad C=-\frac2{n+2}\Sigma,
 \qquad \nabla^SH=\Sigma+\Omega,
$$
where $\Sigma$ and $\Omega$ are the symmetric and skew parts of the
intrinsic derivative of the tangent vector field $H$.  Since each ambient
component of $H$ is a degree-three spherical harmonic,
$$
 \int|\nabla^SH|^2=3n+2.
$$
The divergence-free condition and integration by parts on $S^{n-1}$ give
$$
 \int\inner{\nabla^SH}{(\nabla^SH)^\top}=-(n-2),
$$
and consequently
\begin{equation*}
 \int|\Sigma|^2=n+2,
 \qquad
 \int|\Omega|^2=2n.
\end{equation*}
It follows that
$$
 N=2+\frac4{(n+2)^2}\int|\Sigma|^2
   =\frac{2(n+4)}{n+2},
 \qquad N_1=1.
$$
Furthermore, differentiating $S_H(v)v=H(v)$ yields
$$
 z_{00}=-2H,
 \qquad
 Z_b=\int\left|\nabla^SH+\frac2{n+2}\Sigma\right|^2
 =\frac{3n^2+12n+16}{n+2}.
$$
The remaining entries follow from the block definitions,
\eqref{eq:quartic-WI-common}, and
\eqref{eq:quartic-energy-split}.

Finally, for $S_q$ one has
$$
 a=q,
 \qquad
 b=-\frac1n\nabla^Sq,
 \qquad
 C=\frac1{n(n+1)}
 \bigl(\nabla_S^2q-nq\,g_{S^{n-1}}\bigr).
$$
For a normalized degree-two spherical harmonic,
\begin{equation*}
 \int|\nabla^Sq|^2=2n,
 \qquad
 \int|\nabla_S^2q|^2=2n(n+2),
 \qquad
 \tr(\nabla_S^2q)=-2nq.
\end{equation*}
Thus
$$
 \int|C|^2=\frac{n+4}{n(n+1)},
 \qquad
 N=\frac{(n+2)(n+4)}{n(n+1)},
 \qquad
 N_1=\frac{n+2}{n}.
$$
Differentiating the identity
$S_q(v)v=qv-n^{-1}\nabla^Sq$ gives
$$
 z_{00}=\frac{n+2}{n}\nabla^Sq,
 \qquad
 Z_b=(n+2)^2\int|C|^2
 =\frac{(n+2)^2(n+4)}{n(n+1)}.
$$
Again \eqref{eq:quartic-WI-common} and
\eqref{eq:quartic-energy-split} give all remaining entries.
\end{proof}

\section{Seven-dimensional coefficient }\label{app:seven-dimensional-audit}

\begin{lemma}
\label{lem:seven-contraction}
Let $V$ be a Euclidean $7$-space and let
$$
  u\in \Harm_3(V)\otimes\Lambda^2V^*
$$
satisfy the top normal condition, so that
$$
  g(v):=\iota_vu(v)\in\Harm_2(V)\otimes V^*,
  \qquad \iota_vg(v)=0.
$$
Then
\begin{equation*}
  \|g\|^2\leq\frac23\|u\|^2.
\end{equation*}
\end{lemma}

\begin{proof}
Homogeneously, the normal condition reads
$$
  \iota_xu(x)=|x|^2G(x),
$$
where $G$ is a harmonic homogeneous one-form of degree two satisfying
$\iota_xG(x)=0$.  Define
\begin{equation*}
  u_G(x):=\frac32x^\flat\wedge G(x)-\frac16|x|^2dG(x).
\end{equation*}
Since $G$ is harmonic of degree two,
$$
  \Delta(x^\flat\wedge G)=2dG,
  \qquad
  \Delta(|x|^2dG)=18dG,
$$
so $u_G$ is harmonic.  Cartan's formula for the Euler vector field gives
$\iota_xdG=3G$, hence
$\iota_xu_G(x)=|x|^2G(x)$.  On the unit sphere,
\begin{equation*}
  u_G=v^\flat\wedge g-\frac16d_Sg.
\end{equation*}

Write $G_i(x)=A_{ab;i}x_ax_b$, where
$A_{ab;i}=A_{ba;i}$ and $A_{aa;i}=0$.  Tangency gives
$A_{(ab;i)}=0$.  Contracting one polynomial index with the form index gives
$2A_{ib;i}+A_{ii;b}=0$, and hence $A_{ib;i}=0$.  Thus
$\delta_Sg=0$.  Let $D$ denote the ambient spherical derivative.
Since each ambient component is a degree-two scalar harmonic,
$$
 \int_{S^6}|Dg|^2=14\int_{S^6}|g|^2.
$$
For the tangent field $g$ one has
$D_ag=\nabla^S_ag-\inner{g}{a}v$, and therefore
$|Dg|^2=|\nabla^Sg|^2+|g|^2$.  Hence
$$
  \int_{S^6}|\nabla^Sg|^2=13\int_{S^6}|g|^2.
$$
The coclosed Hodge--Bochner identity on $S^6$ gives
$$
  \int_{S^6}|d_Sg|^2=18\int_{S^6}|g|^2.
$$

Put $u_0:=u-u_G$.  Then $\iota_vu_0=0$.  In components write
$$
 (u_0)_{ij}(x)=U_{abc;ij}x_ax_bx_c,
$$
where $U$ is symmetric in $a,b,c$, skew in $i,j$, and
$U_{aac;ij}=0$.  Exact tangency is the Young relation
$U_{(abc;i)j}=0$.  Contracting $a$ with $i$ in this relation gives
$$
  2U_{ibc;ij}+U_{iic;bj}+U_{ibi;cj}=0,
$$
so $U_{ibc;ij}=0$.  This is precisely $\delta_Su_0=0$.  The radial part of
$u_G$ is pointwise orthogonal to $u_0$, while integration by parts gives
$$
  \int_{S^6}\inner{u_0}{d_Sg}
  =\int_{S^6}\inner{\delta_Su_0}g=0.
$$
Therefore $u_0\perp u_G$, and
$$
  \|u\|^2
  =\|u_0\|^2+\|u_G\|^2
  \geq\|g\|^2+\frac1{36}\|d_Sg\|^2
  =\frac32\|g\|^2.
$$
\end{proof}

\begin{lemma}
\label{lem:seven-degree-two-estimate}
Let $V$ be a Euclidean $7$-space and let
$g\in\Harm_2(V)\otimes V$ satisfy $\inner{g(v)}v=0$.  If
$0\leq\sec_{\cE}\leq1$, then
\begin{equation}\label{eq:seven-degree-two-residual}
  \mathcal W_{\cE}(g)\geq-2\sqrt6\,\|g\|^2,
  \qquad
  \mathcal W_I(g)=13\|g\|^2.
\end{equation}
Consequently, for $\cA=\delta I+(1-\delta)\cE$,
\begin{equation*}
  \mathcal W_{\cA}(g)
  \geq\bigl(13\delta-2\sqrt6(1-\delta)\bigr)\|g\|^2.
\end{equation*}
\end{lemma}

\begin{proof}
The same index contraction as in the preceding proof shows that $g$ is
divergence-free.  At a fixed $v\in S^6$, put
$F(v)=g(v)$ and $\mathsf A=\nabla^SF=\Sigma+\Omega$.  Then
$\tr\Sigma=0$.  The spherical formula for $\mathcal W_{\cE}$ is
\eqref{eq:spherical-w-formula} without its common normalization.  Its
radial term vanishes, and the nonnegative Jacobi terms may be discarded.
The remaining mixed term has the Bianchi decomposition
$$
  \mathfrak m
  =\sum_{a=1}^{6}s_a\cE(F,e_a,v,e_a)
   +\cE(F\wedge v,\omega_\Omega),
$$
where $(s_a)$ are the eigenvalues of $\Sigma$.  Since
$0\preceq\mathsf J_{e_a}^{\cE}\preceq\Id$,
$$
  |\cE(F,e_a,v,e_a)|\leq\frac12|F|.
$$
For a trace-free symmetric endomorphism of a $6$-space,
$\sum_a|s_a|\leq\sqrt6\,|\Sigma|$.  Moreover
$\omega_\Omega$ has at most three simple canonical summands, so
Lemma~\ref{lem:BBK} gives
$$
  |\cE(F\wedge v,\omega_\Omega)|
  \leq\frac23\sqrt3\,|F|\,|\omega_\Omega|.
$$

The scalar components of $F$ have spherical eigenvalue $14$.  If $D$
denotes the ambient spherical derivative, then
$D_aF=\mathsf A(a)-\inner{F}{a}v$, so
$$
 14\|g\|^2=\int|DF|^2=\int|\mathsf A|^2+\|g\|^2.
$$
Thus
$$
  \int|\mathsf A|^2=13\|g\|^2.
$$
The coclosed Hodge--Bochner identity, together with
$dF^\flat=2\omega_\Omega$, gives
$$
  \int|\omega_\Omega|^2=\frac92\|g\|^2,
  \qquad
  \int|\Sigma|^2=4\|g\|^2.
$$
Cauchy--Schwarz bounds each of the two parts of $\mathfrak m$ by
$\sqrt6\|g\|^2$.  This proves the first inequality in
\eqref{eq:seven-degree-two-residual}.  For constant curvature,
$\mathcal W_I=(\lambda_2-1)\|g\|^2=13\|g\|^2$.
\end{proof}

\section{Optimization of the projector threshold}\label{app:projector-optimization}

\begin{remark}[Optimization of the rational constant]
\label{rem:optimized-constant}
The choice $5/13$ is made only to keep
\eqref{eq:quartic-young-one}--\eqref{eq:quartic-gaps} rational.  More
generally, for $a,b>0$ one may use
\begin{align*}
 \frac43\sqrt{A_CZ_C}
 &\leq\frac23\left(aA_C+a^{-1}Z_C\right),\\
 \frac43\sqrt{(n-2)BZ_b}
 &\leq\frac23\left(b(n-2)B+b^{-1}Z_b\right).
\end{align*}
Let $q_*$ be the unique real root of
\begin{equation*}
 648q^3-1404q^2+1026q-251=0.
\end{equation*}
Then
$$
 q_*=0.624487235576400\ldots,
 \qquad
 \delta_*:=\frac{q_*}{1+q_*}
 =0.384421140345138\ldots.
$$
Choose
$$
 a_*=3q_*+\sqrt{9q_*^2-2},
 \qquad
 b_*=\frac{2}{3(2q_*-1)}.
$$
These parameters satisfy
\begin{align*}
 4q_*&=\frac23\left(a_*+\frac2{a_*}\right),\\
 6q_*-1&=\frac{2}{3a_*}
 +\frac23\left(b_*+\frac3{b_*}\right),\\
 2q_*-1&=\frac{2}{3b_*}.
\end{align*}
After substitution in the three rows of
Table \eqref{eq:quartic-table}, the finite-dimensional gaps are
\begin{align*}
 G_{(4,2)}={}&
 \frac{864q_*^2-1350q_*+547}
 {9(18q_*^2-27q_*+11)},\\
 G_{(3,1)}={}&
 \frac{n(324q_*^2-528q_*+239)
       +2808q_*^2-4560q_*+1946}
 {27(n+2)(2q_*-1)(18q_*^2-27q_*+11)},\\
 G_{(2)}={}&
 \frac{n^2(216q_*^2-324q_*+132)
       +n(1404q_*^2-2256q_*+977)
       +2808q_*^2-4560q_*+1946}
 {27n(n+1)(2q_*-1)(18q_*^2-27q_*+11)}.
\end{align*}
Since $0.624<q_*<0.625$, every numerator coefficient and every denominator
in these expressions is positive.  The high-degree calculation remains
positive as well: after writing $n=12+A$, $k=6+B$, the numerator obtained by
completing the square has coefficients
\begin{align*}
 &(18q_*^2+9q_*-4)A^2B+(72q_*^2+18q_*-33)A^2\\
 &+(54q_*^2+9q_*-4)AB^2+(828q_*^2+216q_*-128)AB\\
 &+(2412q_*^2+288q_*-813)A+36q_*^2B^3\\
 &+(1008q_*^2+90q_*-40)B^2
 +(8424q_*^2+1224q_*-880)B\\
 &+19440q_*^2+900q_*-4830,
\end{align*}
all of which are positive on the same interval.  The quadratic-projector
estimate has the smaller threshold $q>\sqrt2/3$.  Consequently, the
proof of Theorem \ref{thm:stable-main} actually yields the conclusion in dimensions
$n\equiv0\pmod4$, $n\geq12$, under strict $\delta_*$-pinching.  Equivalently,
$\delta_*$ is the unique real root in $(0,1)$ of
$$
 3329\delta^3-4209\delta^2+1779\delta-251=0.
$$
\end{remark}

\subsection*{Acknowledgment}
ZSH would like to thank Leiye Xu for helpful discussions.

\subsection*{Declaration of competing interest}
The authors declare that they have no known competing financial interests in this paper.



\subsection*{Data availability}
No data were used for the research described in the article.

\begin{thebibliography}{99}

\bibitem{Berger60}
M.~Berger,
\emph{Pincement riemannien et pincement holomorphe},
Ann. Scuola Norm. Sup. Pisa Cl. Sci. (3) \textbf{14} (1960), no.~2,
151--159.

\bibitem{BettiolGoodman24}
R.~G. Bettiol and M.~J. Goodman,
\emph{Extremality and rigidity for scalar curvature in dimension four},
Selecta Math. \textbf{30} (2024), no.~1, Paper No.~7, 29 pp.

\bibitem{BourguignonKarcher78}
J.-P.~Bourguignon and H.~Karcher,
\emph{Curvature operators: pinching estimates and geometric examples},
Ann. Sci. \'{E}cole Norm. Sup. (4) \textbf{11} (1978), no.~1, 71--92.

\bibitem{Brin75Topology}
M.~Brin,
\emph{The topology of group extensions of C-systems},
Mat. Zametki \textbf{18} (1975), no.~3, 453--465.

\bibitem{Brin75Transitivity}
M.~Brin,
\emph{Topological transitivity of a certain class of dynamical systems, and
flows of frames on manifolds of negative curvature},
Funktsional. Anal. i Prilozhen. \textbf{9} (1975), no.~1, 9--19.

\bibitem{Brin82}
M.~Brin,
\emph{Ergodic theory of frame flows},
in \emph{Ergodic Theory and Dynamical Systems, II} (College Park, 1979/1980),
Progress in Mathematics, vol.~21, Birkh\"auser, Boston, 1982, 163--183.

\bibitem{BrinGromov80}
M.~Brin and M.~Gromov,
\emph{On the ergodicity of frame flows},
Invent. Math. \textbf{60} (1980), 1--7.

\bibitem{BrinKarcher84}
M.~Brin and H.~Karcher,
\emph{Frame flows on manifolds with pinched negative curvature},
Compos. Math. \textbf{52} (1984), no.~3, 275--297.

\bibitem{BurnsPollicott03}
K.~Burns and M.~Pollicott,
\emph{Stable ergodicity and frame flows},
Geom. Dedicata \textbf{98} (2003), 189--210.

\bibitem{CeccheriniScarabottiTolli10}
T.~Ceccherini-Silberstein, F.~Scarabotti and F.~Tolli,
\emph{Representation Theory of the Symmetric Groups:
The Okounkov--Vershik Approach, Character Formulas, and Partition Algebras},
Cambridge Studies in Advanced Mathematics, vol.~121,
Cambridge University Press, Cambridge, 2010.

\bibitem{CekicLefeuvreConnections25}
M.~Ceki\'c and T.~Lefeuvre,
\emph{Isospectral connections, ergodicity of frame flows, and polynomial maps
between spheres},
Ann. Sci. \'{E}cole Norm. Sup. (4) \textbf{58} (2025), no.~1, 203--229.

\bibitem{CekicLefeuvreHolonomy25}
M.~Ceki\'c and T.~Lefeuvre,
\emph{The holonomy inverse problem},
J. Eur. Math. Soc. \textbf{27} (2025), no.~6, 2187--2250.

\bibitem{CekicLefeuvreSemiclassical24}
M.~Ceki\'c and T.~Lefeuvre,
\emph{Semiclassical analysis on principal bundles},
preprint arXiv:2405.14846 (2024).

\bibitem{CLMS22Survey}
M.~Ceki\'c, T.~Lefeuvre, A.~Moroianu and U.~Semmelmann,
\emph{Towards Brin's conjecture on frame flow ergodicity: new progress and
perspectives},
Math. Res. Rep. \textbf{3} (2022), 21--34.

\bibitem{CLMS24}
M.~Ceki\'c, T.~Lefeuvre, A.~Moroianu and U.~Semmelmann,
\emph{On the ergodicity of the frame flow on even-dimensional manifolds},
Invent. Math. \textbf{238} (2024), 1067--1110.

\bibitem{CLMSUnitary24}
M.~Ceki\'c, T.~Lefeuvre, A.~Moroianu and U.~Semmelmann,
\emph{On the ergodicity of unitary frame flows on K\"ahler manifolds},
Ergodic Theory Dynam. Systems \textbf{44} (2024), no.~8, 2143--2172.

\bibitem{CLMSW25}
M.~Ceki\'c, T.~Lefeuvre, A.~Moroianu and U.~Semmelmann,
\emph{Correspondence between Pestov and Weitzenb\"ock identities},
Math. Proc. Cambridge Philos. Soc. \textbf{178} (2025), no.~3, 443--463.

\bibitem{DairbekovSharafutdinov10}
N.~S. Dairbekov and V.~A. Sharafutdinov,
\emph{On conformal Killing symmetric tensor fields on Riemannian manifolds},
Mat. Tr. \textbf{13} (2010), no.~1, 85--145.

\bibitem{DolgopyatKanigowskiRodriguezHertz24}
D.~Dolgopyat, A.~Kanigowski and F.~Rodr\'iguez Hertz,
\emph{Exponential mixing implies Bernoulli},
Ann. of Math. (2) \textbf{199} (2024), no.~3, 1225--1292.

\bibitem{FultonHarris}
W.~Fulton and J.~Harris,
\emph{Representation Theory: A First Course},
Graduate Texts in Mathematics, vol.~129, Springer-Verlag, New York, 1991.

\bibitem{Gray76}
A.~Gray,
\emph{The structure of nearly K\"ahler manifolds},
Math. Ann. \textbf{223} (1976), no.~3, 233--248.

\bibitem{GPSU16}
C.~Guillarmou, G.~P. Paternain, M.~Salo and G.~Uhlmann,
\emph{The X-ray transform for connections in negative curvature},
Comm. Math. Phys. \textbf{343} (2016), no.~1, 83--127.

\bibitem{HasselblattPesin06}
B.~Hasselblatt and Y.~Pesin,
\emph{Partially hyperbolic dynamical systems},
in \emph{Handbook of Dynamical Systems}, vol.~1B,
Elsevier, Amsterdam, 2006, 1--55.

\bibitem{HeilMoroianuSemmelmann16}
K.~Heil, A.~Moroianu and U.~Semmelmann,
\emph{Killing and conformal Killing tensors},
J. Geom. Phys. \textbf{106} (2016), 383--400.

\bibitem{Lefeuvre23}
T.~Lefeuvre,
\emph{Isometric extensions of Anosov flows via microlocal analysis},
Comm. Math. Phys. \textbf{399} (2023), no.~1, 453--479.

\bibitem{Nagy02}
P.-A.~Nagy,
\emph{On nearly-K\"ahler geometry},
Ann. Global Anal. Geom. \textbf{22} (2002), no.~2, 167--178.

\bibitem{PollicottZhang25}
M.~Pollicott and D.~Zhang,
\emph{Rapid mixing for compact group extensions of hyperbolic flows},
Trans. Amer. Math. Soc. \textbf{378} (2025), no.~7, 5011--5056.

\bibitem{Reeder97}
M.~Reeder,
\emph{Exterior powers of the adjoint representation},
Canad. J. Math. \textbf{49} (1997), no.~1, 133--159.

\bibitem{SemmelmannWeingart10}
U.~Semmelmann and G.~Weingart,
\emph{The Weitzenb\"ock machine},
Compos. Math. \textbf{146} (2010), no.~2, 507--540.

\end{thebibliography}

\end{document}

